\documentclass[11pt,a4paper]{article}
\usepackage[margin=23mm]{geometry}
\usepackage[T1]{fontenc}
\usepackage{lmodern,microtype,amsmath,amssymb,amsthm,mathtools,enumitem,xcolor,booktabs}
\usepackage[colorlinks=true,linkcolor=blue!45!black,citecolor=blue!45!black,urlcolor=blue!45!black]{hyperref}
\newtheorem{theorem}{Theorem}[section]
\newtheorem{lemma}[theorem]{Lemma}
\newtheorem{proposition}[theorem]{Proposition}
\newtheorem{corollary}[theorem]{Corollary}
\theoremstyle{remark}\newtheorem{remark}[theorem]{Remark}
\setlength{\parskip}{3pt}
\setlist{nosep}
\allowdisplaybreaks[1]
\title{Thresholds, coherent structures, and dynamics\\of critical focusing dispersive systems\\[5pt]
\large Concentration energy, spectral bifurcation, and transport}
\author{Shuxing Yang}
\date{8 October 2026}
\begin{document}
\maketitle
\begin{abstract}
We connect critical concentration, coherent motion, and conservative
evolution through five analytic results. A bounded modification of a
homogeneous dispersion has an exact concentration energy determined by
the high-frequency optimizer manifold. Strictly below that energy,
critical-mass sublevels are uniformly bounded; a conservative local flow
therefore remains bounded for all time. A concrete anisotropic quartic
class in dimensions one through three includes its local flow proof.
For homogeneous waves, a complementary-resolvent reduction gives the
spectral pair created by small velocity, including the correction due to
translation coupling. Least-mass negative-energy endpoints and regular
zero-energy splitting components force degeneracy of the conserved-charge
response. Internal unitary symmetry gives compactness, quantitative
polarization for arbitrary fields, and profile rigidity uniformly in the
number of components, including a separable Hilbert space. Finally an
exact core cancellation gives a two-sided
localized transport identity for all homogeneous orders $m>1$, including
anisotropic dispersion and critical Hartree interactions. Each theorem
states the operator or flow assumptions it uses.
\end{abstract}

The conditional spectral assertions retain their
nondegeneracy and operator-continuity hypotheses. The new global-flow
theorem is unconditional in the concrete quartic class specified below.


\section*{The central connection}
Critical scaling makes kinetic energy and attraction scale at the same
rate. At critical mass, a concentrating sequence can keep its energy
bounded only by approaching the sharp interpolation optimizer manifold.
A bounded dispersion correction then survives as a finite energy on that
manifold. This identifies the energy at which concentration becomes
possible, and conservation turns the strict inequality below it into a
statement about every time of the actual evolution.

The same coherent structures organize a second question: what becomes
of the critical scaling mode when the wave starts moving? Its interaction
with phase and translation modes determines the spectral growth rate.
The coefficient must be computed after removing translation coupling.
The third question concerns evolution that does not stay near one wave.
A localized dilation identity relates sustained spatial transport to
energy and momentum. These are three distinct mathematical connections;
none is inferred merely from the existence of a stationary profile.

\paragraph{A concrete nonlinear evolution result.}
For $N=1,2,3$, $\varepsilon>0$, consider
\[
 iu_t=k(D)u-|u|^{8/N}u,\qquad
 k(\xi)=|\xi|^4+\varepsilon\frac{|\xi|^{10}}{1+|\xi|^6}.
\]
Let $M_0$ be the pure quartic sharp critical mass and put
$M_\infty=(1+\varepsilon)^{N/4}M_0$.
Every $H^2$ initial datum with
\[
 M(u_0)=M_\infty,\qquad E_k(u_0)<0
\]
produces a global solution with $\sup_{t\in\mathbb R}\|u(t)\|_{H^2}<\infty$.
This class is nonempty: finite dilations of a high-frequency optimizer
have strictly negative energy. The conclusion covers all such initial
data, with no radial symmetry, momentum restriction, or proximity to a
particular wave. It concerns this modified dispersion, whose high-frequency
deficit tends to zero; it is not a blowup theorem for the pure quartic
equation.

\paragraph{Organization and dependencies.}
Section~\ref{sec:r10-foundations} supplies scalar compactness and the
conventions. The concentration-energy, spectral-pair, polarization, and
transport arguments are given in full. The endpoint section applies to
any regular profile at the defined threshold; its finite-splitting
corollary uses the compactification theorem of the preceding companion
manuscript~\cite{companion}. The secondary scale-selection and
minimizer-stability consequences explicitly identify the preceding
variational inputs they import. Their source proofs are included with
this manuscript's source package.

\section{Analytic foundations and conventions}
\label{sec:r10-foundations}
The Fourier transform is unitary, $D=-i\nabla$, and
$P(u)=\operatorname{Im}\int\overline u\nabla u=\int\xi|\widehat u|^2$.
The real Hilbert pairing is the real part of the complex $L^2$ pairing.
Throughout the local-power arguments, $N\ge1$, $m>1$, $s=m/2$,
$r=m/N$, and $p=2+2r$. A principal symbol $a$ is real, even,
$m$-homogeneous, $C^{N+2}$ off zero, and satisfies
$c|\xi|^m\le a(\xi)\le C|\xi|^m$.
Set $A(u)=\int a|\widehat u|^2$, $M(u)=\|u\|_2^2$, and
$E(u)=A(u)/2-\|u\|_p^p/p$. The sharp critical mass $M_c$ is defined by
\begin{equation}\label{eq:r10-basic-GN}
 \|u\|_p^p\le\frac p2\left(\frac{M(u)}{M_c}\right)^r A(u).
\end{equation}
The traveling-wave convention $u(t,x)=e^{i\omega t}U(x-vt)$ gives
$[a(D)+\omega+iv\cdot\nabla]U-|U|^{2r}U=0$.

\begin{lemma}[Normalized optimizer compactness]
The inequality \eqref{eq:r10-basic-GN} is attained. If
$M(f_n)\to M_c$, $A(f_n)\to1$, and $\|f_n\|_p^p\to p/2$, then
some translated subsequence converges strongly in $H^s$ to an optimizer.
The normalized optimizer set is compact modulo translations and phases.
\end{lemma}
\begin{proof}
Sharp Sobolev interpolation gives a finite positive best constant, and
amplitude and spatial scaling normalize a maximizing sequence as stated.
Bounded energy-space sequences that vanish locally in mass converge to
zero in $L^p$: use a bounded-overlap cover by unit balls, local Sobolev
interpolation, and $2<p<2N/(N-m)$ when $N>m$ (all finite $p$ otherwise).
Thus translations give a nonzero weak limit $Q$ and almost-everywhere
convergence. Write $M_1=M(Q)>0$, $A_1=A(Q)>0$, $M_2=M_c-M_1$,
and $A_2=1-A_1$. Quadratic splitting and Brezis--Lieb give
\[
 M_c^r\le M_1^r A_1+M_2^r A_2,\qquad A_1+A_2=1.
\]
If both masses are positive, the right side is strictly below $M_c^r$.
Consequently $M_1=M_c$ and then $A_1=1$. Weak convergence plus
convergence of both quadratic norms gives strong $H^s$ convergence.
The same proof applies to a sequence of optimizers. Constant phase is
proved by the density-norm lemma below; this compactness argument itself
does not need a positivity-preserving kinetic operator.
\end{proof}

\begin{lemma}[Weak-profile regularity and Pohozaev identities]
Every weak $H^s$ local-power profile belongs to $L^\infty\cap H^{m+1}$
and satisfies
\begin{equation}\label{eq:r10-basic-pohozaev}
 v\cdot P=2mE,\qquad \omega M=rA+(2m-p)E.
\end{equation}
No spatial moment is required.
\end{lemma}
\begin{proof}
Choose $\Lambda$ so large that
$a(\xi)-v\cdot\xi+\Lambda\asymp1+|\xi|^m$.
The shifted resolvent maps $L^q$ to $H^{m,q}$, $1<q<\infty$,
by the Mihlin multiplier theorem. For $N>m$, the starting exponent
is $q_0=2N/(N-m)$. Splitting the nonlinearity into $|U|>1$ and
$|U|\le1$, resolvent Sobolev improvement iterates the reciprocal
exponent as $\beta_{j+1}=(1+2r)\beta_j-r$ with
$\beta_0=(1-r)/2$. The ideal iterates
$\beta_j=1/2-(r/2)(1+2r)^j$ cross the boundedness threshold after
finitely many steps. Inputs have exponent greater than one; intermediate
exponents may be weakened to avoid endpoints. The linear contribution
has the stronger improvement $\beta\mapsto\beta-r$ and is handled
separately. When $N=m$ start at any sufficiently large finite exponent;
when $N<m$ the initial space already embeds in $L^\infty$.
Now the nonlinearity is in $L^2$, giving $U\in H^m\subset H^1$.
The $C^1$ map $z\mapsto|z|^{2r}z$, with bounded derivative on the
range of $U$, sends $U$ into $H^1$; one more resolvent application
gives $H^{m+1}$.

Here are the cutoff details for the dilation calculation. Let
$f_R(x)=x\chi(x/R)$ with $\chi$ compactly supported and equal to
one near zero, and let $L_R=f_R\cdot\nabla+\operatorname{div}f_R/2$.
The Fourier kernel of $[a(D),L_R]$ is, up to the transform constant,
\[
 \frac i2(a(\xi)-a(\eta))(\xi+\eta)\cdot\widehat f_R(\xi-\eta).
\]
Splitting comparable and noncomparable frequencies gives
\[
 |a(\xi)-a(\eta)||\xi+\eta|
 \le C|\xi-\eta|(|\xi|+|\eta|)^m.
\]
After division by $\langle\xi\rangle^s\langle\eta\rangle^s$,
the convolution kernel is bounded by
$C|k|\langle k\rangle^s R^{N+1}|\widehat f(Rk)|$, with a
uniform $L^1$ norm for $R\ge1$. Hence the commutators are uniformly
bounded $H^s\to H^{-s}$. Schwartz approximation and the homogeneous
symbol identity give
$\langle U,[a(D),L_R]U\rangle\to mA(U)$.
The drift commutator tends to $-v\cdot P$; its order is one and
$s>1/2$. In the nonlinear pairing integration by parts and dominated
convergence give
\[
 \operatorname{Re}\int |U|^{p-2}\overline U L_RU
 \longrightarrow N(1/2-1/p)\|U\|_p^p=\frac m p\|U\|_p^p.
\]
Testing the profile with $L_RU$ therefore gives $mE=v\cdot P/2$.
Testing with $U$ gives $A+\omega M-v\cdot P-\|U\|_p^p=0$;
combining these two equations proves \eqref{eq:r10-basic-pohozaev}.
\end{proof}

The same regularity proof holds for an individually fixed nonnegative
order-$m$ addition $b$ with the stated Mihlin derivative bounds.
Constants need not be uniform in that addition or in $v$.
For nonhomogeneous minimizers, differentiation of their scalar
mass-preserving dilation energy is legitimate directly by dominated
convergence and minimality; differentiability of the dilation orbit
as an $H^s$-valued map is not being assumed.

\section{Concentration energy and global dynamics at critical mass}
\label{sec:r11-variational}

A high-frequency concentration has a definite limiting energy contribution
from a bounded correction to the homogeneous dispersion.  We identify
that contribution using the full optimizer manifold.  Conservation then
converts the resulting compactness statement into a global theorem for
all initial data below the concentration energy, including data at the
critical mass.  No proximity to a standing wave is needed.

\subsection{The homogeneous problem and the concentration energy}

Fix $N\ge1$, $m>1$, $s=m/2$, $r=m/N$, and $p=2+2r$.
Let $a$ be real, even, positive elliptic, and homogeneous of degree $m$.
Write
\[
 M(f)=\|f\|_2^2,\qquad
 A(f)=\int a(\xi)|\widehat f(\xi)|^2\,d\xi,\qquad
 E_a(f)=\tfrac12 A(f)-\tfrac1p\|f\|_p^p.
\]
Let $M_c$ be the sharp mass in
\begin{equation}\label{eq:r11-sharp-gn}
 \frac1p\|f\|_p^p\le
 \frac12\left(\frac{M(f)}{M_c}\right)^r A(f),
\end{equation}
and define the normalized equality manifold
\[
 \mathcal G=\{Q\in H^s:M(Q)=M_c,\ A(Q)=1,\ E_a(Q)=0\}.
\]
We use the normalized sharp-GN compactness property: whenever
$M(f_n)\to M_c$, $A(f_n)\to1$, and $E_a(f_n)\to0$, there are
translations and a subsequence converging strongly in $H^s$ to an
element of $\mathcal G$.  For the homogeneous symbols considered in
this work this is the established sharp-GN compactness input.  It
requires neither optimizer uniqueness nor nondegeneracy of the real
linearized operator.  In particular $\mathcal G$ is compact modulo
translations; constant phases are already included in this set.

Let $d$ be a bounded, measurable, real, even Fourier multiplier. Assume
that there is a bounded function $d_\infty$ on $S^{N-1}$ such that
\begin{equation}\label{eq:r11-angular-limit}
 \sup_{\theta\in S^{N-1}}|d(R\theta)-d_\infty(\theta)|
       \longrightarrow0\qquad(R\longrightarrow\infty).
\end{equation}
Define
\[
 k=a-d,\qquad
 D(f)=\int d(\xi)|\widehat f(\xi)|^2\,d\xi,\qquad
 E_k(f)=E_a(f)-\tfrac12D(f).
\]
No sign assumption is imposed on $d$ or $d_\infty$.  The quadratic
perturbation is continuous on $L^2$, and $E_k$ is continuous on $H^s$.
Set
\begin{equation}\label{eq:r11-concentration-energy}
 e_\infty=-\frac12\max_{Q\in\mathcal G}
       \int d_\infty(\xi/|\xi|)|\widehat Q(\xi)|^2\,d\xi.
\end{equation}
The maximum exists: its objective is a bounded-multiplier quadratic
form, continuous in $L^2$ and invariant under translations and constant
phase.  Compactness of $\mathcal G$ modulo translations applies.  Since
the limiting multiplier has degree zero, the value does not depend on
the kinetic normalization used for the equality manifold.

\begin{theorem}[Energy of a concentrating sequence]
\label{thm:r11-concentrating-sequence}
Suppose
\[
 M(f_n)\le M_c,\qquad A(f_n)\longrightarrow\infty,\qquad
 \sup_n E_k(f_n)<\infty.
\]
Then $M(f_n)\to M_c$.  Put $t_n=A(f_n)^{1/m}$ and write
$f_n(x)=t_n^{N/2}V_n(t_nx)$.  After extraction and translations,
$V_n\to Q$ strongly in $H^s$ for some $Q\in\mathcal G$, and
\begin{align}
 D(f_n)&\longrightarrow
   \int d_\infty(\xi/|\xi|)|\widehat Q(\xi)|^2\,d\xi,
       \label{eq:r11-deficit-limit}\\
 \liminf_n E_k(f_n)&\ge
   -\frac12\int d_\infty(\xi/|\xi|)|\widehat Q(\xi)|^2\,d\xi
   \ge e_\infty.
       \label{eq:r11-concentrating-liminf}
\end{align}
\end{theorem}

\begin{proof}
The sharp inequality and the mass bound give $E_a(f_n)\ge0$.
Boundedness of $d$ gives
\[
 0\le E_a(f_n)=E_k(f_n)+\tfrac12D(f_n)
      \le \sup_j E_k(f_j)+\tfrac12\|d\|_\infty M_c.
\]
Consequently
\[
 A(V_n)=1,\qquad
 E_a(V_n)=t_n^{-m}E_a(f_n)\longrightarrow0.
\]
Equation~\eqref{eq:r11-sharp-gn} also gives
\[
 0\le\frac12\left[1-\left(\frac{M(f_n)}{M_c}\right)^r\right]
      \le E_a(V_n)\longrightarrow0.
\]
Thus $M(f_n)\to M_c$, and normalized sharp-GN compactness yields
the asserted strong convergence.  Equivalently one may first correct
the masses and kinetic normalizations by factors tending to one and
then undo those normalizations after obtaining a compact subsequence.

Translations do not change Fourier modulus.  Since $d$ is bounded,
strong $L^2$ convergence implies
\[
 \left|\int d(t_n\xi)
      \bigl(|\widehat V_n(\xi)|^2-|\widehat Q(\xi)|^2\bigr)\,d\xi\right|
 \le\|d\|_\infty\|V_n-Q\|_2(\|V_n\|_2+\|Q\|_2)
 \longrightarrow0.
\]
For every $\xi\ne0$, \eqref{eq:r11-angular-limit} gives
$d(t_n\xi)\to d_\infty(\xi/|\xi|)$.  Dominated convergence now
proves~\eqref{eq:r11-deficit-limit}.  Finally
$E_k(f_n)=E_a(f_n)-D(f_n)/2\ge-D(f_n)/2$, proving
\eqref{eq:r11-concentrating-liminf}.
\end{proof}

\begin{corollary}[Bounded energy sublevels and sharpness]
\label{cor:r11-sublevels}
For every $e<e_\infty$,
\begin{equation}\label{eq:r11-uniform-sublevel}
 \sup\{A(f):M(f)\le M_c,\ E_k(f)\le e\}<\infty,
\end{equation}
where the assertion is vacuous if the set is empty.
Conversely, for every $e>e_\infty$ the exact-mass sublevel
$\{f:M(f)=M_c,\ E_k(f)\le e\}$ is unbounded in $H^s$.
\end{corollary}

\begin{proof}
Failure of~\eqref{eq:r11-uniform-sublevel} contradicts
Theorem~\ref{thm:r11-concentrating-sequence}.  For sharpness choose
$Q\in\mathcal G$ attaining the maximum in
\eqref{eq:r11-concentration-energy} and let
$Q_\lambda(x)=\lambda^{N/2}Q(\lambda x)$.  Homogeneity gives
\[
 M(Q_\lambda)=M_c,\qquad A(Q_\lambda)=\lambda^m,\qquad
 E_a(Q_\lambda)=0.
\]
Moreover
\[
 E_k(Q_\lambda)=-\tfrac12\int d(\lambda\xi)|\widehat Q(\xi)|^2\,d\xi
      \longrightarrow e_\infty
      \qquad(\lambda\longrightarrow\infty).
\]
For any fixed $e>e_\infty$ these trials eventually belong to the
stated sublevel while their kinetic energies diverge.
\end{proof}

The strict comparison in this result is optimal for mass--energy
coercivity.  At $e=e_\infty$, boundedness depends on the particular
symbol.  Sharpness of an energy sublevel does not itself construct a
blow-up trajectory.

\begin{remark}[Choice of the zero of frequency]
Replacing $k$ by $k+c$ multiplies solutions by $e^{-ict}$, adds $cM/2$
to their energy, and adds $cM_c/2$ to $e_\infty$.  Thus the difference
$E_k-e_\infty$ at $M=M_c$ is unchanged.  The concentration-energy
comparison does not depend on the arbitrary zero of frequency.
\end{remark}

\subsection{Conservation turns the bound into global dynamics}

For the general-symbol statement, we specify the local evolution input.
Let \textup{(F)} denote the following properties of
\begin{equation}\label{eq:r11-flow}
 i\partial_tu=k(D)u-|u|^{p-2}u:
\end{equation}
for every $u_0\in H^s$ there is a unique maximal solution
$u\in C((T_-,T_+),H^s)$, depending continuously on its initial datum,
conserving $M$ and $E_k$, and satisfying the continuation alternative
that a finite lifespan endpoint forces divergence of the $H^s$ norm.
Ellipticity of an arbitrary symbol is not by itself being used as a
substitute for this local evolution theory.

\begin{theorem}[Global bounded evolution at critical mass]
\label{thm:r11-global-endpoint}
Assume \textup{(F)}.  If
\[
 M(u_0)\le M_c,\qquad E_k(u_0)<e_\infty,
\]
then the solution of~\eqref{eq:r11-flow} is global in both time
directions and
\[
 \sup_{t\in\mathbb R}\|u(t)\|_{H^s}<\infty.
\]
For every strictly subcritical datum $M(u_0)<M_c$, the same conclusion
holds without any restriction on its energy.
\end{theorem}

\begin{proof}
In the first case, conservation places the maximal orbit in the fixed
sublevel of Corollary~\ref{cor:r11-sublevels} with
$e=E_k(u_0)<e_\infty$.  Ellipticity of $a$ and conservation of mass
bound its $H^s$ norm.  The continuation alternative excludes both
finite endpoints.  For the second assertion the direct estimate
\begin{equation}\label{eq:r11-subcritical-coercivity}
 E_k(f)\ge\frac12\left[1-\left(\frac{M(f)}{M_c}\right)^r\right]A(f)
             -\frac12\|d\|_\infty M(f)
\end{equation}
is coercive at each fixed mass below $M_c$; conservation gives the
same conclusion.
\end{proof}

\begin{proposition}[A concrete local theory including anisotropy]
\label{prop:r11-quartic-local}
For $m=4$ and $N\in\{1,2,3\}$, hypothesis \textup{(F)} holds for
every symbol $k=a-d$ in this section.  Consequently
Theorem~\ref{thm:r11-global-endpoint} is unconditional for these
quartic equations, including anisotropic $a$.
\end{proposition}

\begin{proof}
Here $s=2$ and the nonlinear map is $F(z)=|z|^\alpha z$ with
$\alpha=8/N>2$.  Its second derivative, viewing $\mathbb C$ as
$\mathbb R^2$, is locally Lipschitz.  The embeddings
$H^2(\mathbb R^N)\hookrightarrow L^\infty\cap W^{1,4}$ and the
ordinary chain and product rules imply, on every bounded $H^2$ ball,
\begin{equation}\label{eq:r11-local-lipschitz}
 \|F(u)-F(v)\|_{H^2}\le C_R\|u-v\|_{H^2}
 \quad\text{if }\|u\|_{H^2}+\|v\|_{H^2}\le R.
\end{equation}
For example the second-derivative terms have the forms
$F'(u)D^2u$ and $F''(u)(Du)^2$; their differences are bounded using
$H^2\hookrightarrow L^\infty$, $H^2\hookrightarrow W^{1,4}$,
and the local Lipschitz bounds on $F'$ and $F''$.

Since $k$ is real, $e^{-itk(D)}$ is a strongly continuous unitary group
on $H^2$.  Its Duhamel map and~\eqref{eq:r11-local-lipschitz} give a
contraction in $C([-T,T],H^2)$ with $T$ depending only on the initial
$H^2$ norm.  This proves existence, uniqueness, continuous dependence,
and the $H^2$ continuation alternative.

To justify conservation without assuming excessive smoothness of the
non-polynomial power, let $P_R$ be the orthogonal Fourier projection to
$\{\xi:|\xi|\le R\}$ and solve
\[
 i\partial_tu_R=k(D)u_R-P_RF(u_R),\qquad u_R(0)=P_Ru_0.
\]
The band-limited subspace is preserved.  On that subspace $k(D)$ is a
bounded operator on $H^2$, so the solution is differentiable in $H^2$.
The real $L^2$ pairings defining the derivatives of mass and energy are
therefore legitimate.  If $g_R=k(D)u_R-F(u_R)$, then
$\partial_tu_R=-iP_Rg_R$, so
\[
 \frac{d}{dt}E_k(u_R)
   =\operatorname{Re}\langle g_R,-iP_Rg_R\rangle=0.
\]
Likewise $dM(u_R)/dt=0$, since $P_Ru_R=u_R$ and
$\int\overline{u_R}F(u_R)$ is real.  The local contraction bounds are
uniform in $R$.  Strong convergence $P_R\to\mathrm{Id}$ on $H^2$
and the Duhamel equations yield
$u_R\to u$ in $C([-T,T],H^2)$.  Continuity of $M$ and $E_k$ passes
conservation to $u$.  Iterating over the maximal lifespan proves
\textup{(F)}.
\end{proof}

\subsection{Consequences for the two-threshold interpolation}

Let $a_0$ and $a_\infty$ be positive elliptic real-even symbols of degree
$m$, with $h=a_\infty-a_0$ positive elliptic.  Consider
\begin{equation}\label{eq:r11-rational-symbol}
 \begin{gathered}
 k_q(\xi)=a_0(\xi)+h(\xi)\frac{|\xi|^q}{1+|\xi|^q}
         =a_\infty(\xi)-d_q(\xi),\\
 d_q(\xi)=\frac{h(\xi)}{1+|\xi|^q},\qquad q>0.
 \end{gathered}
\end{equation}
Denote the two homogeneous sharp masses by $M_0<M_\infty$.
The following existence and compactness inputs are imported from the
companion study \emph{Thresholds, spectral geometry, and persistent
radiation in mass-critical dispersive equations}, Section~5; they are
not reproved in the present dynamical argument.  That analysis gives
$I(M)=\inf_{M(f)=M}E_{k_q}(f)<0$ with attainment for
$M_0<M<M_\infty$.  At $M=M_\infty$, the infimum is also negative
and attained when $q>m$, and every minimizing sequence is compact in
$H^s$ modulo translations.

\begin{corollary}[All negative-energy endpoint data are bounded]
\label{cor:r11-rational-global}
Assume a conservative local flow with the properties \textup{(F)} for
$k_q$.  For every $q>0$, all initial data of mass less than
$M_\infty$ generate global uniformly $H^s$-bounded solutions.
If $q>m$, the same conclusion holds for every initial datum satisfying
\begin{equation}\label{eq:r11-negative-endpoint-class}
 M(u_0)=M_\infty,\qquad E_{k_q}(u_0)<0.
\end{equation}
For $m=4$, $N=1,2,3$, and $q>4$, the latter assertion requires no
additional local-flow hypothesis.
\end{corollary}

\begin{proof}
For all $q>0$, and every $\eta>0$, homogeneity gives a finite
$C_\eta$ such that
$k_q\ge(1-\eta)a_\infty-C_\eta$.  Hence
\[
 E_{k_q}(f)\ge\frac12
 \left[1-\eta-\left(\frac{M(f)}{M_\infty}\right)^r\right]A_\infty(f)
       -\frac12C_\eta M(f).
\]
For a fixed subcritical mass choose $\eta$ to make the coefficient
positive and use conservation.  If $q>m$, $d_q$ is bounded and tends
uniformly to zero at high frequency.  Thus $d_\infty=0$ and
$e_\infty=0$, so Theorem~\ref{thm:r11-global-endpoint} proves
\eqref{eq:r11-negative-endpoint-class}.  The final assertion follows
from Proposition~\ref{prop:r11-quartic-local}.
\end{proof}

The endpoint result covers arbitrary data in the stated class: it uses
no radial symmetry, momentum restriction, or small distance to a
standing wave.  It excludes both finite-time blow-up and unbounded
infinite-time growth of the energy-space norm.

For $q=m$, the bounded correction has angular limit $h(\theta)$, so
\eqref{eq:r11-concentration-energy} supplies the precise endpoint
comparison.  If $I(M_\infty)<e_\infty$, its strict sublevel is nonempty
and Theorem~\ref{thm:r11-global-endpoint} applies there.  In the
isotropic case $h(\xi)=\varepsilon|\xi|^m$,
\[
 e_\infty=-\tfrac12\varepsilon M_\infty,
 \qquad
 E_{k_m}(f)>-\tfrac12\varepsilon M_\infty
 \quad\text{for every nonzero }f\text{ of mass }M_\infty.
\]
Thus that strict sublevel is empty; no endpoint global conclusion is
being inferred in the isotropic borderline case.  For $q<m$ the
correction is unbounded at high frequency and the concentration theorem
above does not apply.  The variational equality
$I(M_\infty)=-\infty$ in that case does not itself prove finite-time
blow-up.

For $q>m$, if $Q$ is an $a_\infty$ equality profile, then
$E_{k_q}(Q_\lambda)<0$ for every finite $\lambda$, while
$E_{k_q}(Q_\lambda)\to0$ and $A_\infty(Q_\lambda)\to\infty$.
A fixed negative energy therefore gives a bound, but no uniform bound
can cover all negative energies approaching zero.

\subsection{Stability of minimizing sets and the critical mass boundary}

At a mass $c$ where the constrained infimum is attained define
\[
 \Sigma_c=\{U\in H^s:M(U)=c,\ E_k(U)=I(c)\}.
\]
This entire set includes translations and constant phases.  Its
stability makes no uniqueness assertion about individual profiles.

\begin{theorem}[Stability in the correct perturbation class]
\label{thm:r11-set-stability}
Assume \textup{(F)}.
\begin{enumerate}
 \item For the interpolation~\eqref{eq:r11-rational-symbol}, if
 $M_0<c<M_\infty$, then $\Sigma_c$ is orbitally stable under all
 sufficiently small $H^s$ perturbations.
 \item For a bounded correction as above, suppose $c=M_c$,
 $I(c)<e_\infty$, and every minimizing sequence at mass $c$ is
 precompact in $H^s$ modulo translations.  Then $\Sigma_c$ is orbitally
 stable under sufficiently small $H^s$ perturbations whose masses are
 at most $c$.  In particular it is stable on the exact mass surface
 $M=c$.  This applies to~\eqref{eq:r11-rational-symbol} at
 $c=M_\infty$ whenever $q>m$.
\end{enumerate}
In both assertions all sufficiently close initial data in the indicated
perturbation class generate global solutions.  More explicitly, for
any $\epsilon>0$ there is $\delta>0$ such that
$\operatorname{dist}_{H^s}(u_0,\Sigma_c)<\delta$ implies
\[
 \sup_{t\in\mathbb R}\operatorname{dist}_{H^s}(u(t),\Sigma_c)<\epsilon,
\]
with the additional mass restriction in the second assertion.
\end{theorem}

\begin{proof}
The compactness assumption implies that $\Sigma_c$ is bounded in
$H^s$.  Thus closeness to this set makes the mass and energy uniformly
close to $c$ and $I(c)$, respectively.  In the first assertion all
sufficiently close data have mass uniformly below $M_\infty$; the
subcritical coercivity estimate used in
Corollary~\ref{cor:r11-rational-global} makes their solutions global
and bounds their $H^s$ norms uniformly.  In the second assertion choose
$e$ strictly between $I(c)$ and $e_\infty$ and apply
Corollary~\ref{cor:r11-sublevels} to masses at most $c$ and energies
at most $e$.  This gives the same global existence and uniform bound.

If stability failed, there would be initial data $u_{0,n}$ with
$\operatorname{dist}_{H^s}(u_{0,n},\Sigma_c)\to0$ and times $t_n$
such that $f_n=u_n(t_n)$ stays a fixed positive distance from
$\Sigma_c$.  Conservation gives
\[
 M(f_n)\longrightarrow c,\qquad E_k(f_n)\longrightarrow I(c),
\]
and the preceding estimates bound $f_n$ in $H^s$.  Set
$g_n=(c/M(f_n))^{1/2}f_n$.  Then
\[
 M(g_n)=c,\qquad \|g_n-f_n\|_{H^s}\longrightarrow0,\qquad
 E_k(g_n)\longrightarrow I(c).
\]
Compactness of exact-mass minimizing sequences now gives
\[
 \operatorname{dist}_{H^s}(g_n,\Sigma_c)\longrightarrow0.
\]
The same follows for $f_n$, a contradiction.
\end{proof}

The proof is the established variational compactness mechanism for
orbital stability.  The additional concentration-energy bound is what
allows it to operate at the critical mass.  At that boundary the
restriction $M(u_0)\le M_c$ cannot be silently omitted: nearby data can
otherwise cross into the supercritical mass range, where the argument
supplies no uniform kinetic bound.  Neither unconstrained endpoint
stability nor instability is asserted here.

\subsection{Removal of the flat scale direction}

The interpolation also produces an exact measure of the energy cost of
changing the spatial scale of a minimizer.  This statement does not
require spectral nondegeneracy.

\begin{proposition}[Exact scale curvature]
\label{prop:r11-scale-curvature}
Let $U\ne0$ be a constrained energy minimizer for
\eqref{eq:r11-rational-symbol}.  Define
\[
 U_t(x)=e^{Nt/2}U(e^tx),\qquad
 \mathcal K(U)=\left.\frac{d^2}{dt^2}E_{k_q}(U_t)\right|_{t=0}.
\]
Then the scalar energy curve is twice differentiable and, with
$z=|\xi|^q$,
\begin{equation}\label{eq:r11-exact-curvature}
 \mathcal K(U)=\frac q2\int
 h(\xi)\frac{z[(m+q)+(m-q)z]}{(1+z)^3}
                    |\widehat U(\xi)|^2\,d\xi.
\end{equation}
For $0<q\le m$, $\mathcal K(U)>0$.  For $q>m$, minimality still
gives $\mathcal K(U)\ge0$, but positivity does not follow pointwise
from the integrand.
\end{proposition}

\begin{proof}
Let $b(\xi)=h(\xi)|\xi|^q/(1+|\xi|^q)$ and let
$\mathcal D=\xi\cdot\nabla_\xi$.  The derivatives
$\mathcal D^jb$, $j=0,1,2$, are bounded in absolute value by
$C|\xi|^m$.  Therefore differentiating the scalar Fourier energy is
justified by dominated convergence, with no assumption of a spatial
moment or differentiability of $t\mapsto U_t$ in $H^s$.  Homogeneity
gives
\[
 E_{k_q}(U_t)=e^{mt}E_{a_0}(U)
                 +\frac12\int b(e^t\xi)|\widehat U(\xi)|^2\,d\xi.
\]
Its first derivative vanishes at the scalar minimum $t=0$.  Eliminating
$E_{a_0}(U)$ from the second derivative yields
\[
 \mathcal K(U)=\frac12\int(\mathcal D^2-m\mathcal D)b(\xi)
                                    |\widehat U(\xi)|^2\,d\xi.
\]
Since $\mathcal Dh=mh$ and $\mathcal Dz=qz$, direct differentiation
gives
\[
 (\mathcal D^2-m\mathcal D)b
      =q h\frac{z[(m+q)+(m-q)z]}{(1+z)^3},
\]
proving~\eqref{eq:r11-exact-curvature}.  If $q\le m$, its coefficient
is strictly positive for every $\xi\ne0$, and a nonzero $L^2$ function
cannot have all Fourier mass at the origin.  If $q>m$, the coefficient
changes sign at $z=(m+q)/(q-m)$; nonnegativity of the integral follows
instead from minimality.
\end{proof}

For the homogeneous critical optimizer the energy is exactly flat
under mass-preserving dilation.  Proposition~\ref{prop:r11-scale-curvature}
proves that the rational correction removes this degeneracy in the
scale direction when $q\le m$.  It does not identify a complete
Hamiltonian spectral mode.

We give the precise lower-threshold consequence of the selection theorem
established for this interpolation.  Put
\[
 \begin{gathered}
 \alpha=\frac{M}{M_0},\qquad d=\alpha^r-1,\qquad
 U_M(x)=\sqrt\alpha\,\lambda_M^{N/2}V_M(\lambda_Mx),\\
 M(V_M)=M_0,\qquad A_0(V_M)=1.
 \end{gathered}
\]
Let
\[
 \begin{gathered}
 \mathcal G_0=\{Q:M(Q)=M_0,\ A_0(Q)=1,\ E_{a_0}(Q)=0\},\\
 C(Q)=\int h(\xi)|\xi|^q|\widehat Q(\xi)|^2\,d\xi,
 \end{gathered}
\]
and suppose at least one element of $\mathcal G_0$ has finite $C$.
Write $C_{\min}=\min_{\mathcal G_0}C$ and
\[
 \beta_*^q=\frac{m}{(m+q)C_{\min}},\qquad
 K_* =\frac{q}{2(m+q)}\beta_*^m.
\]
The lower-threshold selection theorem of the same companion study
provides the following inputs; the new deduction here is the curvature
asymptotic, not a new proof of profile selection:
\begin{align}
 \lambda_M&\sim\beta_*d^{1/q},&
 I(M)&\sim-\alpha K_*d^{1+m/q},
       \label{eq:r11-selection-recall-a}\\
 \frac{B(U_M)}{A_0(U_M)}&=\frac{m}{m+q}d+o(d),&
 \int\frac{h(\xi)|\xi|^q}{1+|\lambda_M\xi|^q}
                |\widehat V_M(\xi)|^2\,d\xi&\longrightarrow C_{\min},
       \label{eq:r11-selection-recall-b}
\end{align}
where $B(f)=\int b|\widehat f|^2$.  After translations and extraction,
$V_M$ converges strongly in $H^s$ to some $Q\in\mathcal G_0$ with
$C(Q)=C_{\min}$.  The correction-energy convergence in
\eqref{eq:r11-selection-recall-b}, not merely strong $H^s$
convergence, is needed in the next deduction.

\begin{corollary}[Sharp lower-threshold scale curvature]
\label{cor:r11-curvature-asymptotic}
For minimizers satisfying the preceding selection conclusions, for
every $q>0$,
\begin{equation}\label{eq:r11-curvature-asymptotic}
 \mathcal K(U_M)\sim\frac{mq}{2}dA_0(U_M),\qquad
 \frac{\mathcal K(U_M)}{|I(M)|}\longrightarrow m(m+q)
 \qquad(M\downarrow M_0).
\end{equation}
In particular the scale curvature is strictly positive near the lower
threshold also when $q>m$.
\end{corollary}

\begin{proof}
Along a convergent selected-profile subsequence define the nonnegative
finite measures
\[
 d\mu_M(\xi)=
 \frac{h(\xi)|\xi|^q}{1+|\lambda_M\xi|^q}
                       |\widehat V_M(\xi)|^2\,d\xi.
\]
Their total masses tend to $C_{\min}$.  On every compact annulus they
converge to
$h(\xi)|\xi|^q|\widehat Q(\xi)|^2\,d\xi$, whose total mass is also
$C_{\min}$.  Positivity therefore gives tightness away from both zero
and infinity: exhaust the limiting integral by compact annuli and use
total-mass convergence to control the complements.

Formula~\eqref{eq:r11-exact-curvature} becomes
\[
 \frac{\mathcal K(U_M)}{\alpha\lambda_M^{m+q}}
 =\frac q2\int
 \frac{(m+q)+(m-q)|\lambda_M\xi|^q}
      {(1+|\lambda_M\xi|^q)^2}\,d\mu_M(\xi).
\]
The displayed multiplier is uniformly bounded and tends on compact
annuli to $m+q$.  Tightness proves that the right side converges to
$q(m+q)C_{\min}/2$.  Since
$B(U_M)=\alpha\lambda_M^{m+q}\mu_M(\mathbb R^N)$, we have
\[
 \frac{\mathcal K(U_M)}{B(U_M)}\longrightarrow\frac{q(m+q)}2.
\]
Combining this with~\eqref{eq:r11-selection-recall-b} proves the first
assertion.  The energy and scale laws in
\eqref{eq:r11-selection-recall-a}, and
$A_0(U_M)=\alpha\lambda_M^m$, give the second.  Every selected
subsequence has the same scalar limits, so the conclusions hold along
the full family.
\end{proof}

There is also a fixed-scale version.  Let
$b_0(\xi)=a(\xi)|\xi|^q/(1+|\xi|^q)$,
$k_\varepsilon=a+\varepsilon b_0$, and fix $0<\tau<1$.
For the joint limit with
$\alpha_\varepsilon=(1+\varepsilon\tau)^{1/r}$, suppose the
selected minimizers satisfy
$U_\varepsilon/\sqrt{\alpha_\varepsilon}\to Q$ strongly in $H^s$
modulo translations.  Since
$(\mathcal D^2-m\mathcal D)b_0$ is an order-$m$ multiplier,
\begin{equation}\label{eq:r11-joint-curvature}
 \frac{\mathcal K(U_\varepsilon)}
      {\alpha_\varepsilon\varepsilon}
 \longrightarrow\frac12\int
       (\mathcal D^2-m\mathcal D)b_0(\xi)|\widehat Q(\xi)|^2\,d\xi.
\end{equation}
For $q\le m$ the limit is positive.  Thus an arbitrarily weak
perturbation creates a scale curvature of order $\varepsilon$ at the
finite spatial scale selected by a mass excess of the same order.

An oscillation frequency requires a further Hamiltonian spectral
analysis.  The conjugate phase, translation coupling, and interaction
with other spectral modes enter that calculation.  The exact curvature
and its asymptotics above do not presume that such an isolated
oscillatory eigenmode has already been established.

\subsection{Which parts extend to other critical interactions}

The concentration argument uses only a homogeneous energy
$E_a=A/2-\mathcal F$ with $E_a\ge0$ at mass at most $M_c$,
compactness of normalized zero-energy sequences, and a bounded
quadratic correction $-D/2$.  It does not use locality or a polynomial
formula for $\mathcal F$.  Theorems~\ref{thm:r11-concentrating-sequence}
and~\ref{thm:r11-global-endpoint} consequently extend, with the same
proof, to a mass-critical Hartree interaction once its sharp inequality,
optimizer compactness, and conservative local flow have been verified.
The curvature calculation needs only
$\mathcal F(U_t)=e^{mt}\mathcal F(U)$, in addition to the explicit
dispersion.  Set stability further requires compactness of the relevant
constrained minimizing sequences.  This separates the common
concentration mechanism from the interaction-specific analytic inputs.

\section{The spectral pair created by slow translation}
\label{sec:r11-spectral}

The preceding velocity--momentum identities determine the inertia of the
charge Hessian. We now identify the pair of spectral values that emerges
from the critical scaling mode, including its leading coefficient. The
result assumes real-profile nondegeneracy; this hypothesis is not supplied
by phase coercivity alone.

Let $a$ be positive, elliptic, real-even and homogeneous of degree $m>1$.
Consider
\[
 iu_t=a(D)u-|u|^{2m/N}u,
 \qquad
 a(D)U+\omega U+iv\cdot\nabla U-|U|^{2m/N}U=0.
\]
All Hilbert spaces in this section are real unless explicitly complexified.
Let $Q=U_{1,0}$ be a real sharp ground wave and set
\[
 \Lambda Q=\frac N2Q+x\cdot\nabla Q,\qquad
 L_-=a(D)+1-|Q|^{2m/N},\qquad
 L_+=a(D)+1-\left(1+\frac{2m}{N}\right)|Q|^{2m/N}.
\]
We impose the following hypotheses.
\begin{enumerate}
\item The functions $\partial_jQ$ and $\Lambda Q$ have the regularity and
spatial decay needed to differentiate translations and critical scaling
in $L^2$ and in the profile equation. In particular,
$L_+\Lambda Q=-mQ$ holds in the operator sense.
\item $L_-\geq0$, $\ker L_-=\mathbb RQ$, and $L_-^{-1}$ is bounded
on $Q^\perp$. Moreover
$\ker L_+=\operatorname{span}\{\partial_jQ:1\leq j\leq N\}$,
and $L_+$ has a bounded inverse on that kernel's orthogonal complement.
\item There is a $C^1$ local profile branch $U_{\omega,v}$, closed under
critical scaling up to phase and translation. Its reduced action
$d(\omega,v)=E(U_{\omega,v})+\omega M(U_{\omega,v})/2
-v\cdot P(U_{\omega,v})/2$ is $C^2$ and satisfies
$d_\omega=M/2$ and $d_v=-P/2$. Profile differentiation is valid, and in
a translation gauge
\[
 D_{v_j}U_{1,0}=-iL_-^{-1}\partial_jQ.
\]
\item On the real space $\mathcal H=L^2(\mathbb R^N;\mathbb C)$ the
real action Hessians $\mathcal L_v$ are selfadjoint on a common dense
domain. Their kernels are exactly
\[
 K_v=\operatorname{span}_{\mathbb R}
       \{iU_{1,v},\partial_1U_{1,v},\ldots,\partial_NU_{1,v}\}.
\]
These kernel bases vary continuously in $\mathcal H$. The inverse of
$\mathcal L_v$ on $K_v^\perp$, extended by zero on $K_v$, exists and is
norm-continuous as a bounded operator on $\mathcal H$ for small $v$.
\end{enumerate}
The last hypothesis explicitly includes a uniform complementary spectral
gap. Common-domain operator continuity, or suitable norm-resolvent
continuity with an isolated kernel of fixed dimension, provides standard
ways to verify it. It is not being inferred for every dispersion symbol.

Define the real matrix and vector
\[
 T_{ij}=\langle\partial_iQ,L_-^{-1}\partial_jQ\rangle,
 \qquad c_j=\langle\Lambda Q,L_-^{-1}\partial_jQ\rangle,
\]
and the scalar
\begin{equation}\label{eq:r11-spectral-D}
 \mathcal D_Q=
 \langle\Lambda Q,L_-^{-1}\Lambda Q\rangle-c^TT^{-1}c.
\end{equation}
The inverse of $L_-$ always acts on $Q^\perp$.

\begin{theorem}[Small-velocity spectral asymptotics]
\label{thm:r11-spectral-pair}
Under the four hypotheses above, $T$ is positive definite and
$\mathcal D_Q>0$. If $m\ne2$, then for every sufficiently small nonzero
$v$, the nonzero spectrum of
\[
 \mathcal A_v=J\mathcal L_v,\qquad J=-i,
\]
in a fixed neighborhood of zero consists of exactly one simple pair
$\lambda_\pm(v)$, satisfying
\begin{equation}\label{eq:r11-lambda-square}
 \lambda_\pm(v)^2
  =\frac{(m-2)v^TTv}{\mathcal D_Q}+o(|v|^2).
\end{equation}
For $m>2$ this pair is real and
\[
 \lambda_\pm(v)=
 \pm\sqrt{\frac{(m-2)v^TTv}{\mathcal D_Q}}+o(|v|).
\]
In particular, the wave is spectrally unstable. For $1<m<2$ the pair is
purely imaginary, with
\[
 \lambda_\pm(v)=
 \pm i\sqrt{\frac{(2-m)v^TTv}{\mathcal D_Q}}+o(|v|).
\]
The latter assertion concerns the pair near zero and does not, by itself,
exclude unstable spectrum elsewhere.
\end{theorem}

\begin{proof}
Phase coercivity gives $T>0$: a nonzero directional derivative of a
nonzero $L^2$ function cannot vanish. Put
\[
 \Lambda_\perp Q=
 \Lambda Q-\sum_j(T^{-1}c)_j\partial_jQ.
\]
Then
\[
 \mathcal D_Q=
 \langle\Lambda_\perp Q,L_-^{-1}\Lambda_\perp Q\rangle\geq0.
\]
Equality would imply
$(x-\zeta)\cdot\nabla Q+NQ/2=0$ for some $\zeta$.
Thus $Q$ would be fixed by the mass-preserving dilation group about
$\zeta$. A nonzero function homogeneous of degree $-N/2$ cannot belong
to $L^2$: its radial mass integral is proportional to $\int_0^\infty dr/r$.
Hence $\mathcal D_Q>0$. The same Schur-complement formula shows that
$\mathcal D_Q$ is unchanged if a linear combination of translations is
added to $\Lambda Q$, so it is independent of the spatial center.

Write $U_v=U_{1,v}$ and
$g_0=iU_v$, $g_j=\partial_jU_v$. Their span $K_v$ is isotropic:
\begin{equation}\label{eq:r11-isotropic-kernel}
 \langle g_i,Jg_j\rangle=0.
\end{equation}
For the phase--translation pairings this follows from integration of a
mass derivative. For two translations it follows from commutation of
translation charges, or directly from the Fourier representation of
$\langle\partial_iU,-i\partial_jU\rangle_{\mathbb R}$.

Let $\Pi_v$ project orthogonally onto $K_v$, let $\Pi_c=I-\Pi_v$, and
write $\mathcal L_c=\mathcal L_v|_{K_v^\perp}$ and
$J_c=\Pi_cJ\Pi_c$. The operator
\[
 R_v(\lambda)=(\mathcal L_c+\lambda J_c)^{-1}
\]
is extended by zero on $K_v$. A Neumann series and hypothesis 4 show
that this inverse is analytic for $|\lambda|<\rho$, with a common
$\rho>0$ for small $v$. The eigenvalue problem is considered on the
complexification of $\mathcal H$; real pairings below extend bilinearly.

The equation $J\mathcal L_v\psi=\lambda\psi$ is equivalent to
$(\mathcal L_v+\lambda J)\psi=0$. Decompose $\psi=k+w$ with
$k\in K_v$ and $w\in K_v^\perp$. For $\lambda\ne0$ the complementary
equation gives
\[
 w=-\lambda R_v(\lambda)Jk.
\]
By \eqref{eq:r11-isotropic-kernel}, the kernel equation is therefore
exactly $D(\lambda,v)z=0$, where
\begin{equation}\label{eq:r11-reduced-matrix}
 D_{ij}(\lambda,v)
  =\langle g_i,J R_v(\lambda)Jg_j\rangle,
 \qquad k=\sum_{j=0}^Nz_jg_j.
\end{equation}
The common factor suppressed in this equation is $-\lambda^2$.
For $0<|\lambda|<\rho$, analytic block elimination therefore preserves
the eigenvalues and their algebraic multiplicities. Selfadjointness of
$\mathcal L_c$ and skew-adjointness of $J$ imply
\begin{equation}\label{eq:r11-reduced-transpose}
 D(-\lambda,v)=D(\lambda,v)^T.
\end{equation}

Profile differentiation yields
\[
 \mathcal L_v\partial_\omega U=-U,
 \qquad
 \mathcal L_v\partial_{v_j}U=-i\partial_jU=Jg_j.
\]
Components of the parameter derivatives in $K_v$ do not contribute,
by \eqref{eq:r11-isotropic-kernel}. Consequently
\[
 D(0,v)=C d''(1,v)C,
 \qquad C=\operatorname{diag}(1,-I_N).
\]
The lower-right block is $D_{vv}(0,v)=-B(1,v)/2$, where
$B=D_vP$. The first derivative of the branch gives
\[
 P(U_v)=2Tv+o(|v|),\qquad B(1,v)=2T+o(1).
\]
Thus the lower-right block stays invertible for small $\lambda,v$.
Take its scalar Schur complement $f(\lambda,v)$. Identity
\eqref{eq:r11-reduced-transpose} makes $f$ even in $\lambda$.
It is consequently analytic in $z=\lambda^2$ in a common disk, with
coefficients continuous in $v$.

The exact scaling identity for the reduced action gives
\begin{align}
 f(0,v)
 &=\frac{P^TB^{-1}P-(m-1)v\cdot P}{2m^2}\notag\\
 &=\frac{2-m}{m^2}v^TTv+o(|v|^2).
 \label{eq:r11-f-zero}
\end{align}
For completeness, the identity follows by differentiating the scaling
laws for the charges. With $H=d''$ and
$D_*=(m\omega,(m-1)v)^T$, they give
$HD_*=(0,-P/2)^T$ and $H_{vv}=-B/2$. Eliminating the velocity block
shows that its scalar Schur complement is
$[P^TB^{-1}P-(m-1)v\cdot P]/(2m^2\omega^2)$.
This uses no differentiated second-order remainder for the profile.

It remains to compute the coefficient of $\lambda^2$ at $v=0$.
Let $R_0$ be the $L^2$-orthogonal projection of $\Lambda Q/m$ off the
translation kernel, and put $Z_j=L_-^{-1}\partial_jQ$.
At $v=0$ the two real blocks give
\[
 \mathcal L_0=\operatorname{diag}(L_+,L_-),
 \qquad J(h,b)=(b,-h),
\]
\[
 \mathcal L_c^{-1}Jg_0=(-R_0,0),
 \qquad \mathcal L_c^{-1}Jg_j=(0,-Z_j).
\]
Expanding the exact resolvent in \eqref{eq:r11-reduced-matrix}, and
writing $b_j=\langle R_0,Z_j\rangle$, gives
\begin{align*}
 D_{00}(\lambda,0)
  &=\lambda^2\langle R_0,L_-^{-1}R_0\rangle+O(\lambda^4),\\
 D_{0v}(\lambda,0)&=\lambda b^T+O(\lambda^2),\\
 D_{v0}(\lambda,0)&=-\lambda b+O(\lambda^2),\\
 D_{vv}(\lambda,0)&=-T+O(\lambda).
\end{align*}
The scalar Schur complement is even. Its coefficient is therefore
\[
 \partial_z f(0,0)=
 \langle R_0,L_-^{-1}R_0\rangle-b^TT^{-1}b
 =\frac{\mathcal D_Q}{m^2}>0.
\]
The final equality uses invariance of the Schur expression under
addition of translations to the scaling derivative. Analyticity with
a common radius and continuity in $v$ now give
\[
 f(\lambda,v)=\frac{2-m}{m^2}v^TTv+o(|v|^2)
 +\lambda^2\left(\frac{\mathcal D_Q}{m^2}+o(1)\right)
 +O(\lambda^4),
\]
where the last remainder is uniform for small $v$.

The analytic scalar equation, regarded as an equation in $z=\lambda^2$,
has one simple real root near zero, by
$\partial_zf(0,0)>0$ and uniform analytic inversion. That root obeys
\eqref{eq:r11-lambda-square}. Its sign is the sign of $m-2$, since
$T>0$, and it is nonzero for nonzero sufficiently small $v$.
Its two square roots are simple because
$\partial_\lambda f=2\lambda\partial_zf\ne0$ there.
The invertible complementary and velocity blocks exclude further
nonzero spectrum in a sufficiently small fixed neighborhood.
\end{proof}

\begin{remark}[Center, frequency, and linear time scale]
If $Q(-x)=Q(x)$, then $c=0$ by parity. Without this symmetry, omitting
$c^TT^{-1}c$ gives an incorrect coefficient in general.
Critical scaling also gives the exact spectral covariance
\[
 \lambda(\omega,v)=
 \omega\lambda\bigl(1,v\omega^{-(m-1)/m}\bigr).
\]
Thus for $m>2$ and fixed $\omega>0$,
\[
 \lambda_+(\omega,v)=\omega^{1/m}
 \sqrt{\frac{(m-2)v^TTv}{\mathcal D_Q}}+o(|v|).
\]
The reciprocal growth rate is of order $|v|^{-1}$. This is a statement
about linear exponential growth, not an established nonlinear escape
time. When $m=2$, the displayed leading coefficient vanishes; a general
homogeneous degree-two symbol need not be a quadratic form, so an exact
Galilean degeneracy is not asserted.
\end{remark}

\begin{corollary}[Critical Hartree interactions]
Theorem~\ref{thm:r11-spectral-pair} holds, with the same coefficient,
for the critical Hartree flow
\[
 iu_t=a(D)u-(K*|u|^2)u,
 \qquad K(x)=|x|^{-m},\quad 1<m<N,
\]
whenever its real ground wave and local branch satisfy the same four
hypotheses with
\[
 L_-=a(D)+1-(K*Q^2),\qquad
 L_+h=L_-h-2Q\bigl(K*(Qh)\bigr).
\]
It also holds for a positive-type anisotropic interaction of the same
homogeneity under those operator and regularity assumptions.
\end{corollary}
\begin{proof}
The proof above uses the nonlinearity through phase and translation
invariance, conjugation symmetry, critical scaling, and the stated
real-block operators. These properties hold for the Hartree energy.
In particular $L_+\Lambda Q=-mQ$, the parameter equations and the
charge-scaling identities are unchanged. No step in the resolvent
reduction uses locality of $L_+$. Thus the proof applies verbatim.
\end{proof}

\begin{remark}[Relation to established instability theory]
The general mechanism of a spectral pair leaving zero at charge
Hessian degeneracy belongs to established Hamiltonian bifurcation
theory. Here the additional calculation is its coefficient for
mass-critical translation branches, including the coupling to the
translation modes. The theorem proves spectral instability for
$m>2$ under its explicit assumptions. Nonlinear orbital instability
requires a sufficiently differentiable local flow and the associated
linear evolution bounds or unstable-manifold construction in a
specified phase space; it is not inferred here for every admissible
symbol and dimension. Likewise, real-profile nondegeneracy is not
proved by this section.
\end{remark}

\section{Least traveling mass and degenerate neutral modes}
\label{sec:r11-endpoint}

Let $m>1$, $s=m/2$, and let $a$ be a real-even, positive elliptic
$m$-homogeneous symbol with the regularity needed for the energy-space
profile equation and its Pohozaev identities.  On
$X=H^s(\mathbb R^N;\mathbb C)$, considered as a real space, put
\begin{gather*}
 E(U)=\frac12\int a(\xi)|\widehat U(\xi)|^2\,d\xi
       -\frac1p\int|U|^p,\qquad p=2+\frac{2m}{N},\\
 M(U)=\|U\|_2^2,\qquad
 P(U)=\operatorname{Im}\int\overline U\nabla U.
\end{gather*}
An exact profile satisfies
\begin{equation}\label{eq:r11-end-profile}
 F(U,\omega,v):=a(D)U+\omega U+iv\cdot\nabla U-|U|^{p-2}U=0.
\end{equation}
The action, its real Hessian, and the Hamiltonian structure are
\[
 S_{\omega,v}=E+\omega M/2-v\cdot P/2,
 \qquad L=S''_{\omega,v}(U),\qquad J=-i.
\]
Assume the negative-energy profile set is nonempty and define
\[
 M_{\mathrm{tr}}=\inf\{M(U):F(U,\omega,v)=0
                  \text{ for some }(\omega,v),\ E(U)<0\}.
\]
The identity $v\cdot P=2mE$ is understood with this normalization.
All parameter families below are defined on open parameter sets; they
are $C^1$ into $X$, differentiation of the profile equation is valid in
$X^*$, and their reduced action is $C^2$.

\subsection{Mass extrema and the reduced action}

For a full profile family $U_\theta$, $\theta=(\omega,v)$, write
\[
 d(\theta)=S_\theta(U_\theta),\qquad H=d''(\theta),
 \qquad B=D_vP(U_\theta),\qquad e_\omega=(1,0,\ldots,0)^T.
\]
Stationarity implies the exact identities
\begin{equation}\label{eq:r11-end-legendre}
 \nabla d=(M/2,-P/2)^T,\qquad
 E=d-\theta\cdot\nabla d,\qquad
 \nabla E=-H\theta,\qquad \nabla M=2He_\omega.
\end{equation}
In particular, $H\theta=0$ is a null-vector identity and does not
mean that the entire matrix $H$ vanishes.

If the charges of the family obey critical scaling covariance,
differentiating that covariance and using the symmetry of $H$ gives
\begin{equation}\label{eq:r11-end-scaling}
 m\omega M_\omega+(m-1)v\cdot M_v=0,\qquad
 m\omega P_\omega+(m-1)Bv=P,\qquad M_v=-P_\omega.
\end{equation}
This covariance holds for a family locally unique modulo phase and
translation, because critical dilation produces another exact profile
with parameters $(\lambda^m\omega,\lambda^{m-1}v)$.

\begin{theorem}[An attained negative-energy endpoint]
\label{thm:r11-end-attained}
Suppose a member of a full profile family satisfies
$E(U)<0$ and $M(U)=M_{\mathrm{tr}}$.  Then
\begin{equation}\label{eq:r11-end-attained}
 M_\omega=0,\qquad M_v=0,\qquad P_\omega=0,
 \qquad He_\omega=0.
\end{equation}
If the family has critical scaling covariance, then additionally
\begin{equation}\label{eq:r11-end-response}
 (m-1)Bv=P,\qquad v^TBv=\frac{v\cdot P}{m-1}<0.
\end{equation}
Thus the charge Hessian is singular, whether or not $B$ is invertible.
\end{theorem}
\begin{proof}
Negative energy persists in a parameter neighborhood.  By definition of
$M_{\mathrm{tr}}$, mass has a local minimum there at the profile under
consideration.  Its first derivatives vanish.  Equations
\eqref{eq:r11-end-legendre} and \eqref{eq:r11-end-scaling} give the
conclusions, and the strict sign follows from $v\cdot P=2mE<0$.
\end{proof}

The supplied-family hypothesis can be replaced by the following standard
local profile assumptions.  Suppose $F$ is $C^1$ as a map into $X^*$,
$L:X\to X^*$ is selfadjoint Fredholm, and
\begin{equation}\label{eq:r11-end-kernel}
 \ker L=K:=\operatorname{span}_{\mathbb R}
             \{iU,\partial_1U,\ldots,\partial_NU\}.
\end{equation}
Assume the profile and projected solutions have enough regularity for
these generators to belong to $X$ and for the symmetry Noether pairings
below to hold.  These are local regularity assumptions, not an assertion
that translations are differentiable at every point of $X$.

\begin{lemma}[Local profile family from a symmetry slice]
\label{lem:r11-end-slice}
Under these assumptions, the profile has a full $C^1$ local parameter
family, unique up to phase and translation, and its charges have critical
scaling covariance.
\end{lemma}
\begin{proof}
The generators in \eqref{eq:r11-end-kernel} are independent: a real
dependence implies a constant-coefficient first-order equation for $U$,
forcing its $L^2$ Fourier transform onto an affine hyperplane, or forcing
$U=0$.  Take the fixed $L^2$-orthogonal slice to $K$ in $X$.  By
selfadjoint Fredholmness, the restriction of $L$ is an isomorphism from
that slice onto the annihilator of $K$ in $X^*$.  Apply the implicit-function
theorem to the projected equation.  Its residual is a linear combination
of the original generators, identified with elements of $X^*$.
Gauge and translation invariance of $S_{\omega,v}$ give
\[
 \langle F(U_{\rm new},\omega,v),iU_{\rm new}\rangle_{\mathbb R}=0,
 \qquad
 \langle F(U_{\rm new},\omega,v),\partial_jU_{\rm new}\rangle_{\mathbb R}=0.
\]
The pairing matrix of the original and new generators is invertible near
$U$.  Thus all residual coefficients vanish.  Local slice uniqueness
gives uniqueness modulo symmetries.  Applying that uniqueness to the
critically dilated profiles proves covariance of their charges.
\end{proof}

For example, a positive free spectral gap
$\omega>\sup_\xi(v\cdot\xi-a(\xi))$, together with compactness of the
nonlinear potential perturbation, supplies the Fredholm hypothesis.
The kernel and local regularity assumptions must still be verified.

\subsection{A degenerate charge response yields secular linear growth}

\begin{theorem}[An extended neutral Jordan chain]
\label{thm:r11-end-jordan}
Assume \eqref{eq:r11-end-kernel} and the regularity of
Lemma~\ref{lem:r11-end-slice}.  If $0\ne q=(q_\omega,q_v)\in\ker H$,
set
\[
 Z=D_\theta U[q],\qquad
 g=q_\omega iU-q_v\cdot\nabla U.
\]
There exists $W\in X$ such that
\begin{equation}\label{eq:r11-end-chain}
 LW=iZ,\qquad JLW=Z,\qquad JLZ=g,\qquad JLg=0.
\end{equation}
The chain has length three.  If $JL$ generates the linearized evolution
on a phase space containing these vectors in its generator domain, then
\begin{equation}\label{eq:r11-end-polynomial}
 e^{tJL}W=W+tZ+\frac{t^2}{2}g.
\end{equation}
In the quotient by the phase and translation tangent space $K$, this
solution grows at least linearly as $t\to+\infty$.
\end{theorem}
\begin{proof}
Parameter differentiation gives
$LU_\omega=-U$ and $LU_{v_j}=-i\partial_jU$, so $JLZ=g$.
Solvability of $LW=iZ$ is equivalent to $iZ$ annihilating $K$.
Its pairings with $iU$ and $\partial_jU$ are, respectively,
$DM(U)[Z]/2$ and $DP_j(U)[Z]/2$.  These vanish precisely because
$Hq=0$.  Selfadjoint Fredholm solvability gives $W$.
Independence of the symmetry generators implies $g\ne0$; hence
$Z\notin K$.  The finite Jordan chain proves
\eqref{eq:r11-end-polynomial}.  In $X/K$ its value is $[W]+t[Z]$,
whose norm is at least $t\|[Z]\|-\|[W]\|$.
\end{proof}

At the attained endpoint of Theorem~\ref{thm:r11-end-attained}, one may
take $q=e_\omega$.  The conclusion is secular growth of the linearized
equation modulo symmetries.  It asserts neither a positive-real-part
eigenvalue nor nonlinear orbital instability.

\subsection{Zero-energy components of a threshold limit}

The preceding compactification result for $M_{\mathrm{tr}}$ says that a
mass-minimizing negative-energy sequence, normalized to kinetic energy
one, either converges strongly modulo translation to a negative-energy
profile or splits strongly into finitely many zero-energy profiles with
common parameters.  In the splitting alternative,
\[
 \sum_jM(V_j)=M_{\mathrm{tr}},\qquad
 \omega_*=(m/N)/M_{\mathrm{tr}}>0.
\]
The following statements add restrictions to those alternatives.

\begin{theorem}[Degeneracy of each genuinely split component]
\label{thm:r11-end-split}
Let $V$ be a zero-energy exact profile with $M(V)<M_{\mathrm{tr}}$
belonging to a full profile family.  Then
\begin{equation}\label{eq:r11-end-zero}
 \nabla_\theta E=0,\qquad H\theta=0.
\end{equation}
For a scaling-covariant family with $\omega\ne0$, one also has
\begin{equation}\label{eq:r11-end-zero-response}
 Bv=-P,\qquad M_v=-P/\omega,\qquad M_\omega=0.
\end{equation}
Consequently every component in a split threshold limit with at least two
nonzero components has singular charge Hessian whenever it has a regular
profile slice as above; it then has the neutral chain of
Theorem~\ref{thm:r11-end-jordan}.
\end{theorem}
\begin{proof}
All sufficiently nearby members of the family still have mass less than
$M_{\mathrm{tr}}$.  Their energy is nonnegative by definition of that
threshold.  Thus $E$ has a local minimum at $V$ and
\eqref{eq:r11-end-legendre} gives \eqref{eq:r11-end-zero}.
Differentiating $E=v\cdot P/(2m)$ gives $Bv=-P$.  Substitution into
\eqref{eq:r11-end-scaling} gives $M_v=-P/\omega$, and its first
identity, together with $v\cdot P=0$, gives $M_\omega=0$.
In a genuine finite split every component has mass strictly below the
sum, and $\theta_*\ne0$ because $\omega_*>0$.
\end{proof}

If the family is rotationally equivariant for an isotropic equation,
$P=p(c)v/c$ for $c=|v|>0$.  At a component in
Theorem~\ref{thm:r11-end-split}, zero energy gives $p(c)=0$ and
$Bv=-P$ gives $p'(c)=0$.  Hence $B=0$; when $\omega\ne0$,
\eqref{eq:r11-end-zero-response} shows that the \emph{entire} matrix
$H$ vanishes.  Isotropy of the equation alone does not assert rotational
equivariance of every profile family.

\begin{proposition}[A moving zero-energy endpoint]
\label{prop:r11-end-moving}
Suppose $E(U)=0$, $M(U)\le M_{\mathrm{tr}}$, $v\ne0$, and $U$
belongs to a full profile family.  Then $H$ is singular.
\end{proposition}
\begin{proof}
If $H$ were nonsingular, the map $\theta\mapsto(M,P)$ would be a local
$C^1$ diffeomorphism.  Along exact profiles, stationarity gives
\[
 dE=-\frac\omega2\,dM+\frac v2\cdot dP.
\]
Choose $dM<0$ and then choose $dP$ with a sufficiently negative component
along $v$ that $dE<0$.  A small positive perturbation in this charge
direction has both negative energy and mass below $M_{\mathrm{tr}}$,
a contradiction.
\end{proof}

\begin{remark}[The standing-wave exception]
At a normalized zero-energy endpoint $\omega>0$.  If $v=0$ and $P=0$,
scaling gives $H(m\omega,0)^T=(0,-P/2)^T=0$, so $H$ is again singular.
The remaining case $v=0$, $P\ne0$ is not excluded.  The present argument
does not prove that every complex stationary profile is real up to a
constant phase, or that every such profile has zero momentum.  Thus the
correct statement for a single regular zero-energy endpoint retains the
possibility of a standing profile with nonzero momentum.
\end{remark}

\subsection{The exact turning-point slope in an isotropic branch}

\begin{proposition}\label{prop:r11-end-turning}
Suppose an isotropic rotationally equivariant branch, at fixed
$\omega>0$, reaches an attained negative-energy mass threshold at speed
$c_*>0$.  Write $P=p(c)v/c$, and assume the differentiability used below.
Then
\[
 p(c_*)<0,\qquad
 p'(c_*)=\frac{p(c_*)}{(m-1)c_*}<0.
\]
The matrix $B$ is negative definite at $c_*$, and $H$ has $N$ positive
eigenvalues and a one-dimensional kernel.  If the mass minimum satisfies
$M''(c_*)>0$, define
\[
 \kappa(c)=\frac{p(c)^2}{p'(c)}-(m-1)cp(c).
\]
Then
\begin{equation}\label{eq:r11-end-kappaprime}
 \kappa(c_*)=0,\qquad
 \kappa'(c_*)=-m\omega(m-1)c_*M''(c_*)<0.
\end{equation}
Near $c_*$ the positive index of $H$ drops from $N+1$ to $N$ as $c$
increases through $c_*$.
\end{proposition}
\begin{proof}
Use \eqref{eq:r11-end-response}; the transverse eigenvalues of $B$
are $p(c_*)/c_*<0$, and the longitudinal one is $p'(c_*)<0$.
At the endpoint the first row and column of $H$ vanish and
$H_{vv}=-B/2$, proving its stated inertia.
Stationarity of $S$ and $E=cp/(2m)$ give
\[
 m\omega M'=(m-1)cp'-p.
\]
Differentiate this identity once.  Direct differentiation of $\kappa$
at $c_*$, using $(m-1)c_*p'=p$, gives
\[
 \kappa'=(2-m)p-(m-1)^2c_*^2p''
          =-m\omega(m-1)c_*M''.
\]
Finally the Schur complement of the $v$ block in $H$ equals
$\kappa/(2m^2\omega^2)$, obtained from
\eqref{eq:r11-end-scaling} and $H_{vv}=-B/2$.  The $v$ block remains
positive definite nearby and $\kappa$ changes from positive to negative.
\end{proof}

The relation between conserved-charge degeneracy and extended neutral
Jordan chains belongs to classical Hamiltonian solitary-wave theory,
including the work of Grillakis--Shatah--Strauss and
Comech--Pelinovsky.  The conclusions here specify what that mechanism
forces at the least-negative-energy traveling-mass threshold and in its
finite splitting.  The turning formula alone is not an eigenvalue
bifurcation theorem; that requires the additional spectral hypotheses
and resolvent expansion stated separately.

\section{Polarization rigidity with arbitrarily many internal components}
\label{sec:r11-systems}

The density geometry has a consequence which is stronger than a
componentwise extension: near criticality, a common-frequency traveling
wave cannot acquire independent spatial profiles in its internal
components. The number of components may be unbounded.

\subsection{Class of systems and main statements}
Let $\mathcal H$ be any nonzero separable complex Hilbert space.
Finite-dimensional $\mathbb C^d$ and $\ell^2$ are both allowed.
Fix $N\ge1$, $m>1$, $s=m/2$, and a real-even, positive elliptic,
$m$-homogeneous symbol $a$, of class $C^{N+2}$ away from zero.
For $U\in H^s(\mathbb R^N;\mathcal H)$ put
\[
 \rho_U=\|U\|_{\mathcal H}^2,\quad
 M(U)=\int\rho_U,\quad
 A(U)=\int a(\xi)\|\widehat U(\xi)\|_{\mathcal H}^2\,d\xi.
\]
We use either of the following interactions:
\begin{align*}
 &\kappa=1+m/N,
 &&\mathcal F(U)=\int\rho_U^\kappa,
 &&V_\rho=\rho^{\kappa-1};\\
 &\kappa=2,
 &&\mathcal F(U)=\mathcal D(\rho_U,\rho_U),
 &&V_\rho=K*\rho.
\end{align*}
In the second case $1<m<N$, the real even kernel $K$ is homogeneous
of degree $-m$, its interaction form $\mathcal D$ is positive definite,
and
\[
 |K(x)|+|x||\nabla K(x)|\le C|x|^{-m}.
\]
As before, positivity of the Fourier transform of $K$ almost everywhere
suffices for positive definiteness. Pointwise positivity is not needed.
The scalar sharp inequality defines $M_c$ by
\begin{equation}\label{eq:r11-vectorGN}
 \mathcal F(u)\le\kappa
       \left(\frac{M(u)}{M_c}\right)^{\kappa-1}A(u).
\end{equation}

Allow a common scalar addition $b\ge0$, real and even, such that
\begin{equation}\label{eq:r11-vectorcone}
 \xi\cdot\nabla b\le mb,\qquad
 |\partial^\beta b(\xi)|\le C_\beta|\xi|^{-|\beta|}(1+|\xi|^m),
 \quad |\beta|\le N+2.
\end{equation}
The constants need not be bounded over the class. Write
\[
 B_b(U)=\int b(\xi)\|\widehat U(\xi)\|^2\,d\xi,
 \qquad E_b(U)=\frac{A(U)+B_b(U)}2-\frac{\mathcal F(U)}{2\kappa}.
\]
A common-frequency traveling profile satisfies
\begin{equation}\label{eq:r11-vectorprofile}
 [a(D)+b(D)+\omega+iv\cdot\nabla]U-V_{\rho_U}U=0.
\end{equation}
All components have the same $a,b,\omega,v$. This is the profile
of $u(t,x)=e^{i\omega t}U(x-vt)$.

\begin{theorem}[Sharp interpolation and dimension-independent compactness]
\label{thm:r11-vectorcompact}
Inequality \eqref{eq:r11-vectorGN} has the same sharp constant for
$\mathcal H$-valued functions. Every optimizer has the form
\[
 U(x)=\zeta Q(x),\qquad \|\zeta\|_{\mathcal H}=1,
\]
where $Q$ is a real scalar optimizer, up to a constant phase.
If $M(U_n)\to M_c$, $A(U_n)\to1$, and
$\mathcal F(U_n)\to\kappa$, then the sequence is strongly compact
in $H^s$ after translations and internal unitaries. This remains
true when the internal dimensions vary: embed all spaces in
$\ell^2$ before applying the unitaries.
\end{theorem}

\begin{theorem}[Nonpositive-energy rigidity, uniform in internal dimension]
\label{thm:r11-vectorgap}
There is $\delta>0$, depending only on $N,m,a$ and the scalar
interaction, such that every nonzero weak profile
\eqref{eq:r11-vectorprofile} satisfying
\[
 M(U)<M_c+\delta,\qquad E_b(U)\le0
\]
has
\[
 E_b(U)=0,\qquad v=0,\qquad D_b(U)=0,\qquad U=\zeta q,
\]
where
\[
 D_b(U)=\int[mb-\xi\cdot\nabla b]\|\widehat U\|^2
\]
and $q$ is a real scalar stationary profile for $a+b$ and $\omega$.
Necessarily $M(U)\ge M_c$. In particular, every negative-energy
profile has mass at least $M_c+\delta$, uniformly over $\mathcal H$,
the cone \eqref{eq:r11-vectorcone}, all frequencies and all velocities.
If $mb-\xi\cdot\nabla b>0$ almost everywhere away from zero, no
nonzero nonpositive-energy profile lies in this mass window.
\end{theorem}

\begin{corollary}[Exact polarization near criticality]
\label{cor:r11-vectorpolarization}
For $b=0$, there is $\eta>0$, independent of $\mathcal H$, such that
every nonzero weak profile satisfying
\[
 \left|\frac{M(U)}{M_c}-1\right|+\frac{|E_0(U)|}{A(U)}<\eta
\]
is exactly polarized: $U=\zeta q$ for a scalar traveling profile $q$.
Positive energy is allowed in this statement.
\end{corollary}

\subsection{Sharp interpolation and complete polarization}
The local interaction is the $\kappa$-th power of the $L^\kappa$
density norm. The Hartree interaction is the square of its positive
definite Hilbert density norm. Thus both have the form
$\mathcal F(U)=\|\rho_U\|_{\mathcal B}^\kappa$ for a strictly convex
real density space $\mathcal B$.
In an internal orthonormal basis $U=(u_j)$, the triangle inequality,
scalar sharp interpolation, and sequence H\"older give
\begin{align*}
 \mathcal F(U)^{1/\kappa}
 &\le\sum_j\mathcal F(u_j)^{1/\kappa}\\
 &\le\kappa^{1/\kappa}M_c^{-(\kappa-1)/\kappa}
   \sum_j A(u_j)^{1/\kappa}M(u_j)^{(\kappa-1)/\kappa}\\
 &\le\kappa^{1/\kappa}M_c^{-(\kappa-1)/\kappa}
              A(U)^{1/\kappa}M(U)^{(\kappa-1)/\kappa}.
\end{align*}
Finite partial sums converge in the density norm because the same
bound applied to their tails tends to zero. A scalar optimizer placed
in one component proves sharpness.

Split the components further into real scalar functions $f_j$.
Equality forces every active $f_j$ to be a scalar optimizer, all
$f_j^2$ to be proportional, and all kinetic-to-mass ratios to agree.
The same conclusions hold after any real orthogonal rotation of a
pair $f_i,f_j$. If $f_j/f_i$ took both signs of a nonzero constant,
the squared sum and difference could not be proportional; when the
constant has magnitude one their nonzero supports would be disjoint.
Thus each $f_j$ is one real constant multiple of a fixed scalar
function. This proves complete polarization, including for countably
many components.

\subsection{Compactness modulo internal unitaries}
Ordinary local compactness does not hold directly in an infinite
internal dimension: a fixed spatial shape can move through mutually
orthogonal internal vectors. The following argument removes precisely
that loss of compactness.

\begin{lemma}[Internal-unitary local compactness]
\label{lem:r11-internalcompact}
Let $X=H^s(\mathbb R^N;\mathbb C)$ with $s>0$. From every bounded
sequence in $X\widehat\otimes\ell^2$, one can extract a subsequence
and internal unitaries $W_n$ such that $W_nU_n$ converges weakly in
$X\widehat\otimes\ell^2$ and strongly in
$L^2_{\mathrm{loc}}(\mathbb R^N;\ell^2)$.
\end{lemma}
\begin{proof}
Fix an orthonormal basis $(\phi_k)$ of the scalar Hilbert space $X$.
Write its internal coefficient vectors as
$c_{nk}=\langle\phi_k,U_n\rangle_X\in\ell^2$.
Gram--Schmidt the first $n$ vectors and extend the resulting finite
orthonormal list to an internal orthonormal basis. A unitary $W_n$
can therefore be chosen so that
\[
 W_nc_{nk}\in\operatorname{span}\{e_1,\ldots,e_k\},\qquad k\le n.
\]
For each fixed $k$, these vectors have a convergent subsequence in a
fixed finite-dimensional space. Diagonal extraction gives strong
convergence of every coefficient vector. Weak compactness gives a
weak limit of $W_nU_n$, with exactly those coefficients.

For $\chi\in C_c^\infty$, the scalar multiplication map
$K_\chi:X\to L^2$, $f\mapsto\chi f$, is compact. If $P_L$ is the
scalar projection onto the first $L$ basis vectors, then
\[
 \|K_\chi(I-P_L)\|_{X\to L^2}\longrightarrow0.
\]
Tensoring with the identity on $\ell^2$ preserves this norm.
The transformed $P_LU_n$ converge strongly because there are only
$L$ coefficient vectors; after multiplication by $\chi$, the tails
are uniformly small. This proves local strong convergence.
\end{proof}

\begin{proof}[Compactness assertion of Theorem~\ref{thm:r11-vectorcompact}]
The localized Sobolev inequality has constants independent of internal
dimension: use the Hilbert-valued Bessel-potential representation and
its norm triangle inequality, then spatial localization and the sum of
component quadratic norms. Thus local mass vanishing for an
$H^s$-bounded sequence implies vanishing in $L^p(\mathcal H)$ for
$2<p<2N/(N-2s)$, or every finite $p>2$ if $2s\ge N$.
In the local case take $p=2\kappa$; in the Hartree case use
$p=4N/(2N-m)$ and Hardy--Littlewood--Sobolev. Hence
$\mathcal F(U_n)\to\kappa>0$ excludes local mass vanishing.

Translate a nonzero amount of local mass into a fixed ball and use
Lemma~\ref{lem:r11-internalcompact}. We obtain a nonzero weak limit
$U$, strong local $L^2$ convergence, and by interpolation strong local
convergence in the interaction exponent. Put $R_n=U_n-U$.
The quadratic quantities split. For local powers the Brezis--Lieb
argument for $\|U\|_{\mathcal H}^{2\kappa}$ gives
\[
 \mathcal F(U_n)=\mathcal F(U)+\mathcal F(R_n)+o(1).
\]
For Hartree, the term $\operatorname{Re}\langle U,R_n\rangle$
tends to zero in $L^{p/2}$ by local convergence and the small
$L^p$ tail of $U$. Also
$\mathcal D(\rho_U,\rho_{R_n})\to0$: approximate $U$ in
$L^p(\mathcal H)$ by a compactly supported smooth finite-dimensional
field. Its potential is bounded and decays at infinity, while $R_n$
tends locally to zero in $L^2$. HLS controls the approximation error.
This proves the same splitting for Hartree.

Let $M_1>0,M_2$ and $A_1,A_2$ be the limiting contributions.
Then $M_1+M_2=M_c$, $A_1+A_2=1$, and sharp interpolation yields
\[
 M_c^{\kappa-1}\le
 M_1^{\kappa-1}A_1+M_2^{\kappa-1}A_2.
\]
Both masses cannot be positive, since each would be strictly below
$M_c$. Thus $M_2=0$, and the inequality then forces $A_2=0$.
The convergence is strong in $H^s$. The limit is an optimizer,
already proved to be completely polarized.
\end{proof}

\subsection{Phase coercivity and the internal symmetry kernel}
For completeness, the scalar phase fact used below follows directly
from the density norm. Normalize a real optimizer by $M(Q)=M_c$,
$A(Q)=1$, and put $\omega_*=(\kappa-1)/M_c$. The form
\[
 L_Q=a(D)+\omega_*-V_{Q^2}
\]
is nonnegative: differentiate the sharp inequality at
$Q+i\sqrt t\,h$. If $h$ belongs to its kernel, write
$x=A(h)$ and $y=M(h)/M_c$. The supporting-functional inequality
for the density norm, sharp interpolation for $h$, and weighted
arithmetic--geometric mean give
\[
 x+(\kappa-1)y
 \le\kappa x^{1/\kappa}y^{(\kappa-1)/\kappa}
 \le x+(\kappa-1)y.
\]
Equality implies $x=y$ and $h^2=yQ^2$ by strict convexity.
The kernel of a nonnegative form is linear; apply the same conclusion
to $h+Q$ to get $h=cQ$. The potential is compact as a form
$H^s\to H^{-s}$, by local Rellich and H\"older in the local case,
or the compact-potential approximation and HLS in the Hartree case.
Consequently
\begin{equation}\label{eq:r11-vectorcoercivity}
 \langle h,L_Qh\rangle\ge c_Q\|h\|_{H^s}^2,
 \qquad \langle h,Q\rangle_{L^2}=0.
\end{equation}
This holds for complex $h$ by real--imaginary splitting, and it can
be summed over arbitrarily many components with the same constant.

At $Qe_1$, write a perturbation as $(h+ik)e_1+Z$ with $Z\perp e_1$.
The Hessian of $E+\omega_*M/2$ is
\[
 \langle h,L_+h\rangle+\langle k,L_Qk\rangle
 +\sum_{j\ge2}\bigl(
  \langle\operatorname{Re}Z_j,L_Q\operatorname{Re}Z_j\rangle
 +\langle\operatorname{Im}Z_j,L_Q\operatorname{Im}Z_j\rangle\bigr).
\]
Its kernel in the phase and internal transverse sector is exactly
\[
 \{Q\eta:\ \eta\in\mathcal H,
                 \operatorname{Re}\langle\eta,e_1\rangle=0\}.
\]
For $\mathcal H=\mathbb C^d$, this is the $(2d-1)$-dimensional real
tangent space to the internal unitary orbit. Any additional full
Hessian kernel lies in the scalar real sector $L_+$; no assumption
on that sector is made.

\subsection{Quantitative polarization without real-sector nondegeneracy}
The preceding compactness and phase coercivity also give a quantitative
statement for arbitrary fields, not just stationary or traveling ones.
Define the scale-invariant interpolation deficit
\[
 \mathfrak d(U)=1-
 \frac{\mathcal F(U)}{\kappa(M(U)/M_c)^{\kappa-1}A(U)}
\]
and the squared distance to a completely polarized field by
\[
 \mathfrak p(U)=\inf_{\substack{\|\zeta\|_{\mathcal H}=1\\q\in H^s}}
 \left\{\frac{M(U-\zeta q)}{M(U)}
             +\frac{A(U-\zeta q)}{A(U)}\right\}.
\]
Both are defined for every nonzero $U\in H^s(\mathcal H)$.

\begin{theorem}[Dimension-independent quantitative polarization]
\label{thm:r11-quantitativepolarization}
There is $C<\infty$, depending only on the fixed scalar dispersion
and interaction, such that
\begin{equation}\label{eq:r11-polarizationdeficit}
 \mathfrak p(U)\le C\mathfrak d(U)
\end{equation}
for every nonzero Hilbert-valued field $U$ and every internal dimension.
No nondegeneracy or uniqueness of the scalar real-sector ground state
is required.
\end{theorem}
\begin{proof}
Amplitude and spatial scaling preserve both quantities, so impose
$M(U)=M_c$, $A(U)=1$. Write
$\|U\|_X^2=M(U)/M_c+A(U)$.
Since $\mathfrak p(U)\le2$, failure of a uniform constant would
produce $U_n$ with $\mathfrak d(U_n)\to0$ and
$\mathfrak d(U_n)/\mathfrak p(U_n)\to0$.
After the symmetries of Theorem~\ref{thm:r11-vectorcompact},
$U_n\to Qe_1$ strongly in $H^s$.
Align $\int QU_n$ with the positive real first internal direction
and write $U_n=(u_n,Z_n)$. Each component of $Z_n$ is complex
$L^2$-orthogonal to $Q$, and $u_n\to Q$, $Z_n\to0$ in $H^s$.

The interaction expansion in these transverse directions is
\begin{equation}\label{eq:r11-transverseexpansion}
 \mathcal F(u_n,Z_n)=\mathcal F(u_n)
 +\kappa\langle V_{|u_n|^2},\|Z_n\|^2\rangle
 +o(\|Z_n\|_X^2).
\end{equation}
For local powers, set $a=|u_n|^2$ and $b=\|Z_n\|^2$ pointwise.
When $1<\kappa\le2$, the nonnegative Taylor remainder of
$(a+b)^\kappa$ after its linear term in $b$ is at most $C_\kappa b^\kappa$.
When $\kappa>2$, it is at most
$C_\kappa(a^{\kappa-2}b^2+b^\kappa)$.
H\"older and the dimension-independent $L^{2\kappa}$ Sobolev bound
make these respectively $O(\|Z_n\|_X^{2\kappa})$ and
$O(\|Z_n\|_X^4+\|Z_n\|_X^{2\kappa})$, since $u_n$ stays bounded.
Both are $o(\|Z_n\|_X^2)$.
For Hartree the remainder is exactly
$\mathcal D(\|Z_n\|^2,\|Z_n\|^2)=O(\|Z_n\|_X^4)$.

Put $m_n=M(Z_n)$ and $a_n=A(Z_n)$. Scalar sharp interpolation gives
\[
 \mathcal F(u_n)\le
 \kappa(1-m_n/M_c)^{\kappa-1}(1-a_n).
\]
Combining its expansion with \eqref{eq:r11-transverseexpansion}
and potential form convergence yields
\begin{align*}
 \kappa-\mathcal F(U_n)
 &\ge\kappa\left[
 a_n+\frac{\kappa-1}{M_c}m_n
       -\langle V_{Q^2},\|Z_n\|^2\rangle\right]
       -o(\|Z_n\|_X^2)\\
 &=\kappa\sum_j\langle Z_{nj},L_QZ_{nj}\rangle
       -o(\|Z_n\|_X^2)\\
 &\ge c\|Z_n\|_X^2.
\end{align*}
The constant is independent of component number. Taking the candidate
$e_1u_n$ in the distance gives
$\mathfrak p(U_n)\le\|Z_n\|_X^2$.
As $\mathfrak d(U_n)=(\kappa-\mathcal F(U_n))/\kappa$,
this contradicts the assumed ratio and proves the theorem.
\end{proof}

The theorem controls internal polarization rather than the full
scalar shape. It therefore remains valid even when scalar optimizer
uniqueness or real-sector quantitative stability is unavailable.

\subsection{Profile identities and nonpositive-energy rigidity}
Weak profiles have $U\in L^\infty(\mathcal H)\cap H^{m+1}(\mathcal H)$.
Here the regularity argument is unaffected by the internal dimension.
Choose a shifted scalar resolvent with symbol comparable to
$1+|\xi|^m$. Its $L^p(\mathcal H)\to H^{m,p}(\mathcal H)$ bounds
hold by Hilbert-valued Mihlin estimates. In the local case, for
$N>m$, the reciprocal Sobolev exponents improve under
$\beta\mapsto(1+2m/N)\beta-m/N$, starting at
$\beta=(1-m/N)/2$; weakening intermediate exponents avoids endpoints
and gives boundedness in finitely many steps. The cases $N\le m$
start at arbitrarily large finite exponent or at boundedness.
For Hartree, the near kernel is in $L^j$ for $j<N/m$ and the far
kernel is bounded. Since $\rho_U\in L^{N/(N-m)}$, the near potential
belongs to arbitrarily large finite $L^t$. Resolvent iteration with
$t>N/m$ again gives boundedness. Bounded nonlinearity coefficients
then give $H^m\subset H^1$; differentiating the nonlinearity once
gives $H^{m+1}$. The derivative of the local vector map is bounded
on bounded sets. In the Hartree case use
$\nabla\rho_U\in L^1\cap L^2$ and the near/far split.

Let $P(U)=\int\xi\|\widehat U\|^2$, and abbreviate $A,B,M,E,D_b$
by $A,B,M,E,D$. Cutoff dilation and testing by $U$ give
\begin{equation}\label{eq:r11-vectorpohoz}
 v\cdot P=2mE-D,\qquad
 \omega M=(\kappa-1)(A+B)+(2m-2\kappa)E-D.
\end{equation}
These require no spatial moment. For the kinetic cutoff use
$f_R=x\chi(x/R)$ and the uniformly bounded commutator with
$f_R\cdot\nabla+\operatorname{div}f_R/2$.
The symbol difference estimate
\[
 |(a+b)(\xi)-(a+b)(\eta)||\xi+\eta|
 \le C|\xi-\eta|[1+(|\xi|+|\eta|)^m]
\]
gives a uniform $H^s\to H^{-s}$ form bound after dividing by the
two Sobolev weights. Schwartz approximation gives the limit
$\int\xi\cdot\nabla(a+b)\|\widehat U\|^2$.
The local nonlinear term is integrated by parts. The Hartree term
is symmetrized into
\[
 -\frac14\iint\rho_U(x)\rho_U(y)
    [f_R(x)-f_R(y)]\cdot\nabla K(x-y)\,dx\,dy,
\]
which tends to $m\mathcal F(U)/4$ by HLS and the uniform Lipschitz
bound on $f_R$. Finite internal projection followed by convergence
justifies all component sums.

Write $U=X+iY$ in an internal basis and set
$\widetilde\omega=\omega+D/M$. The imaginary equation and
$P(U)=-2\int\langle Y,\nabla X\rangle$ give
\begin{align}
 mE={}&\langle Y,[a(D)+\widetilde\omega-V_{\rho_U}]Y\rangle
       +B_b(Y)+D\left(\frac12-\frac{M(Y)}M\right),
       \label{eq:r11-vectorshift}\\
 \widetilde\omega M={}&(\kappa-1)(A+B)+(2m-2\kappa)E.
       \label{eq:r11-vectorshiftomega}
\end{align}

\begin{proof}[Proof of Theorem~\ref{thm:r11-vectorgap}]
Sharp interpolation excludes $M<M_c$ for $E\le0$.
Suppose a sequence violating the stated rigidity had $M_n\to M_c$.
Critical scaling preserves the entire scalar cone, so set $A_n=1$.
Then
\[
 0\le B_n\le(M_n/M_c)^{\kappa-1}-1\longrightarrow0,
 \qquad E_n\to0,\qquad\mathcal F(U_n)\to\kappa.
\]
Theorem~\ref{thm:r11-vectorcompact} gives $U_n\to Qe_1$ strongly
in $H^s$ after translations and internal unitaries.
A further unitary aligns the internal vector
$\int Q(x)U_n(x)\,dx$ with the positive real first coordinate,
without losing this convergence. Thus every imaginary component
of $Y_n$ is orthogonal to $Q$.

Equation~\eqref{eq:r11-vectorshiftomega} gives
$\widetilde\omega_n\to\omega_*$, even if $D_n$ is unbounded.
Potential form continuity and \eqref{eq:r11-vectorcoercivity} give
\[
 \langle Y_n,[a(D)+\widetilde\omega_n-V_{\rho_{U_n}}]Y_n\rangle
 \ge c\|Y_n\|_{H^s(\mathcal H)}^2.
\]
For large $n$, $M(Y_n)/M_n<1/4$. All terms on the right of
\eqref{eq:r11-vectorshift} are nonnegative. Since $E_n\le0$,
we obtain $E_n=0$, $Y_n=0$ and $D_n=0$.
The imaginary equation reduces to $v_n\cdot\nabla X_n=0$.
A nonzero $L^2$ function cannot be constant in a nonzero direction,
so $v_n=0$.

Every real component belongs to the kernel of
$T_n=a(D)+b_n(D)+\omega_n-V_{\rho_{U_n}}$.
Since $D_n=0$, $\omega_n=\widetilde\omega_n$. This form is strictly
positive on $Q^\perp$, by phase coercivity and $b_n\ge0$.
Its real kernel has dimension at most one, since two independent
kernel vectors would have a nonzero combination in $Q^\perp$.
Therefore all components are multiples of one real scalar stationary
profile, the desired rigidity. The sequential contradiction proves a
uniform $\delta$.
\end{proof}

\subsection{Polarization of low positive-energy traveling waves}
For pure $a$, a common-frequency profile also satisfies the phase-boost
identity
\begin{equation}\label{eq:r11-vectorvelocity}
 vM(U)=\int\nabla a(\xi)\|\widehat U(\xi)\|^2\,d\xi.
\end{equation}
To justify it without moments, choose real $\chi\in C_c^\infty$
equal to $x_j$ near zero, put $\chi_R=R\chi(x/R)$, and vary the
profile action along $e^{i\theta\chi_R}U$. The density is unchanged.
The commutator kernel contains
$(a(\xi)-a(\eta))R^{N+1}\widehat\chi(R(\xi-\eta))$.
The bound
$|a(\xi)-a(\eta)|\le C|\xi-\eta|(|\xi|+|\eta|)^{m-1}$
yields a uniform form bound on $H^s$.
For Schwartz fields the limit is Fourier differentiation; density
and the uniform bound extend it to the energy space, giving
\eqref{eq:r11-vectorvelocity} componentwise and after summation.

\begin{proof}[Proof of Corollary~\ref{cor:r11-vectorpolarization}]
If no uniform $\eta$ existed, take unpolarized profiles with
$M_n\to M_c$, $A_n=1$ and $E_n\to0$.
Then $\mathcal F(U_n)\to\kappa$, so $U_n\to Qe_1$ strongly
in $H^s$ after symmetries. The right side of
\eqref{eq:r11-vectorvelocity} converges to zero: $Q$ is real,
$|\widehat Q|^2$ is even, and $\nabla a$ is odd of order $m-1$.
Thus $v_n\to0$.
Equation~\eqref{eq:r11-vectorpohoz} gives $\omega_n\to\omega_*$.
The complex self-adjoint frozen profile form
\[
 T_n=a(D)+\omega_n+iv_n\cdot\nabla-V_{\rho_{U_n}}
\]
converges in form norm to $L_Q$; the drift has order one, strictly
below $m$. Hence it is positive on the complex codimension-one space
$Q^\perp$. Its complex kernel has dimension at most one. All internal
components of $U_n$ belong to this kernel, so they are constant complex
multiples of a single scalar profile, a contradiction.
\end{proof}

\paragraph{Scope of the internal symmetry.}
A non-scalar matrix symbol or an internal rotation
$u(t,x)=e^{it\Omega}U(x-vt)$ with a non-scalar Hermitian $\Omega$
is a different problem. In the latter case different components may
occupy different spectral values of a common frozen spatial operator.
The common-frequency kernel argument does not apply. The statements
above also do not classify arbitrary time-dependent vector solutions.

\section{Homogeneous transport, anisotropic dispersion, and nonlocal attraction}
\label{sec:r11-transport}

The transport argument extends to general homogeneous symbols without a
sign assumption on the localized principal commutator. Its key is an
exact cancellation on a spatially supported core. We obtain a common
time-dependent result for local critical powers and Hartree attraction.
The quantity localized in the general theorem is the full homogeneous
kinetic density; the stronger gradient conclusion for isotropic quartic
dispersion is considered separately at the end.

\subsection{Equations and solution class}
Fix $m>1$, $N\ge1$, and $s=m/2$. Let $a$ be real, even, homogeneous
of degree $m$, $C^{N+2}$ away from zero, and positive elliptic:
\[
 c|\xi|^m\le a(\xi)\le C|\xi|^m.
\]
Set
\[
 A(u)=\int a(\xi)|\widehat u(\xi)|^2\,d\xi,
 \quad M(u)=\|u\|_2^2,
 \quad P(u)=\int\xi|\widehat u(\xi)|^2\,d\xi.
\]
Momentum is finite since $m>1$. We treat either
\begin{align}
 iu_t&=a(D)u-|u|^{p-2}u,
 &p&=2+2m/N,
 &E(u)&=\tfrac12A(u)-\tfrac1p\|u\|_p^p,
 \label{r11tr:local}
\end{align}
or, for $1<m<N$,
\begin{align}
 iu_t&=a(D)u-(K*|u|^2)u,
 &E(u)&=\tfrac12A(u)-\tfrac14H(u),
 \label{r11tr:hartree}\\
 H(u)&=\iint K(x-y)|u(x)|^2|u(y)|^2\,dx\,dy.
 \notag
\end{align}
Here $K$ is real, even, $C^1$ off zero, homogeneous of degree $-m$, and
\begin{equation}\label{r11tr:kernel}
 |K(x)|+|x||\nabla K(x)|\le C_K|x|^{-m}.
\end{equation}
Positive type is unnecessary for the transport theorem: all nonlinear
error estimates use the absolute Riesz interaction. This differs from
the positivity needed for the phase-kernel argument.

Throughout this section, $u\in C([0,\infty);H^s)$ is a strong energy
solution with conserved $M,P,E$, satisfying
\begin{equation}\label{r11tr:bound}
 \sup_{t\ge0}\|u(t)\|_{H^s}\le B<\infty.
\end{equation}
We assume the solution is obtained, on compact time intervals, as a
strong energy-space limit of sufficiently smooth solutions, so that
the integrated localized virial identity follows from the estimates
below. No local well-posedness theorem for arbitrary anisotropic
symbols is asserted here.

Choose $\chi\in C_c^\infty$ with $0\le\chi\le1$, equal to one on
$|x|\le1$ and zero on $|x|\ge2$, and put $\eta=1-\chi$.
For a center $c(t)$ define
\begin{align}
 T_R^c(t)&=\bigl\||D|^s[\eta((x-c(t))/R)u(t)]\bigr\|_2^2,
 \label{r11tr:cut-tail}\\
 D_R^c(t)&=\int_{|x-c(t)|>R}\bigl||D|^s u(t,x)\bigr|^2\,dx.
 \label{r11tr:literal-tail}
\end{align}
These are distinct quantities, but the cutoff commutator below makes
their localization properties equivalent on a bounded energy ball.

\subsection{The kinetic commutator is exact on the core}
Choose a real smooth compactly supported vector field $f$ equal to $x$
on $|x|\le4$. Set
\[
 f_R(x)=Rf(x/R),\qquad
 L_R=f_R\cdot\nabla+\tfrac12\operatorname{div}f_R,
 \qquad B_R=-iL_R,
 \qquad C_R=[a(D),L_R].
\]
Thus $L_R$ is skew-adjoint, while $B_R$ and $C_R$ are symmetric forms.
Write
\[
 Q_R(w)^2=\||D|^s w\|_2^2+R^{-m}\|w\|_2^2.
\]

\begin{lemma}[Uniform form bound and core cancellation]
\label{lem:r11tr-core}
For all $R>0$ and $z,w\in H^s$,
\begin{equation}\label{r11tr:form}
 |\langle z,C_Rw\rangle|+
 |\langle z,ma(D)w\rangle|\le C Q_R(z)Q_R(w).
\end{equation}
If $z$ is supported in $|x|\le2R$, then
\begin{equation}\label{r11tr:exact-core}
 \langle z,[C_R-ma(D)]z\rangle=0.
\end{equation}
Consequently, on a bounded $H^s$ ball and for $R\ge1$,
\begin{equation}\label{r11tr:kinetic-tail}
 \left|\langle u,C_Ru\rangle-mA(u)\right|
 \le C_B\left(\||D|^s(\eta_Ru)\|_2+R^{-s}\right).
\end{equation}
All statements also hold about any fixed center.
\end{lemma}
\begin{proof}
At $R=1$ the Fourier kernel of $C_1$, up to a fixed normalization and
the factor $i$, is
\[
 \tfrac12[a(\xi)-a(\zeta)](\xi+\zeta)\cdot
                   \widehat f(\xi-\zeta).
\]
The homogeneous derivative bound gives
\[
 |a(\xi)-a(\zeta)|\,|\xi+\zeta|
 \le C|\xi-\zeta|(|\xi|+|\zeta|)^m.
\]
After division by $\langle\xi\rangle^s\langle\zeta\rangle^s$,
the right side is at most
$C|\xi-\zeta|\langle\xi-\zeta\rangle^s$. Its product with
$|\widehat f(\xi-\zeta)|$ is integrable. Schur's test therefore
bounds $C_1:H^s\to H^{-s}$. Under the unitary rescaling
$v(y)=R^{N/2}w(Ry)$, $C_R$ becomes $R^{-m}C_1$.
This proves \eqref{r11tr:form}; the term $ma(D)$ is immediate.

For smooth compactly supported $z$ in the stated ball,
$L_Rz=(x\cdot\nabla+N/2)z$. Homogeneity and differentiation of
$A(\lambda^{N/2}z(\lambda\cdot))=\lambda^mA(z)$ give
\[
 2\operatorname{Re}\langle a(D)z,L_Rz\rangle=mA(z),
\]
which is \eqref{r11tr:exact-core}. Approximation within a slightly
larger ball and \eqref{r11tr:form} extend the equality to supported
$H^s$ functions. This is equality of quadratic forms, not a claim
that the defect operator annihilates $z$ pointwise.

Let $z=\chi_Ru$, $w=\eta_Ru$ and $\mathcal D_R=C_R-ma(D)$.
The core equality gives
\[
 \langle u,\mathcal D_Ru\rangle
 =\langle w,\mathcal D_Rw\rangle
       +2\operatorname{Re}\langle z,\mathcal D_Rw\rangle.
\]
The smooth-cutoff product estimate
\[
 \||D|^s(\chi_Ru)\|_2
 \le C\bigl(\||D|^su\|_2+R^{-s}\|u\|_2\bigr)
\]
and \eqref{r11tr:form} now prove \eqref{r11tr:kinetic-tail}.
The product estimate itself follows at scale one from the convolution
inequality $|\xi|^s\le C_s(|\zeta|^s+|\xi-\zeta|^s)$ and Young's
inequality, followed by rescaling.
\end{proof}

No favorable sign of the principal symbol
$\nabla a(\xi)\cdot Df_R(x)\xi$ is used. In particular, the proof
does not infer such a sign from ellipticity or convexity of $a$.

\begin{lemma}[Cutoff energy and the literal kinetic tail]
\label{lem:r11tr-tail-equivalence}
Put $\delta=\min(s,1)$. On the ball \eqref{r11tr:bound},
\begin{equation}\label{r11tr:tail-equivalence}
 (T_R^c)^{1/2}\le(D_R^c)^{1/2}+C_BR^{-\delta},\qquad
 (D_{2R}^c)^{1/2}\le(T_R^c)^{1/2}+C_BR^{-\delta}.
\end{equation}
\end{lemma}
\begin{proof}
For $0<s<1$, the Fourier inequality
$||\xi|^s-|\zeta|^s|\le|\xi-\zeta|^s$ and Young's inequality give
\[
 \|[|D|^s,\eta_R]u\|_2\le CR^{-s}\|u\|_2.
\]
For $s\ge1$, use
\[
 ||\xi|^s-|\zeta|^s|
 \le C_s|\xi-\zeta|
             (|\zeta|^{s-1}+|\xi-\zeta|^{s-1})
\]
to obtain
\[
 \|[|D|^s,\eta_R]u\|_2
 \le C\left(R^{-1}\||D|^{s-1}u\|_2+R^{-s}\|u\|_2\right).
\]
Only the compact part $-\chi_R$ contributes to the commutator.
The uniform bound is $C_BR^{-\delta}$. Since $\eta_R=0$ inside
the radius-$R$ ball and $\eta_R=1$ outside the radius-$2R$ ball,
the triangle inequality proves both assertions.
\end{proof}

\subsection{The common time-dependent estimate}
Let $c$ be Lipschitz with $\|c'\|_\infty\le L$ and set
\begin{equation}\label{r11tr:virial}
 J_R(t)=2\langle u(t),B_R^{c(t)}u(t)\rangle
 =2\operatorname{Im}\int f_R(x-c(t))\cdot\overline u\nabla u.
\end{equation}
For $s<1$, the last expression denotes the symmetric $H^{1/2}$
quadratic form, defined by density.

\begin{lemma}[Two-sided localized virial estimate]
\label{lem:r11tr-virial}
Set
\[
 \gamma=\min\left(\tfrac12,\tfrac1m\right),\qquad
 \beta=\min\left(\tfrac m2,1\right),\qquad
 \widetilde\beta=\min\left(\tfrac m2,\tfrac2m\right).
\]
There are constants $C_0,C_1$ depending on the fixed equation,
$M,B,L$ and cutoffs such that, for $R\ge1$,
\begin{align}
 |J_R(t)|&\le C_0R,\label{r11tr:endpoint}\\
 |J_R'-4mE+2c'\cdot P|
 &\le C_1\left((T_R^c)^\gamma+R^{-\beta}\right),
 \label{r11tr:virial-cut}\\
 |J_R'-4mE+2c'\cdot P|
 &\le C_1\left((D_R^c)^\gamma+R^{-\widetilde\beta}\right).
 \label{r11tr:virial-literal}
\end{align}
The derivative assertions are understood as integrated identities with
an error bounded by the integral of the displayed right side.
\end{lemma}
\begin{proof}
The scale-one first-order Fourier-Schur form estimate and rescaling give
\[
 |\langle u,B_Ru\rangle|
 \le C\bigl(R\||D|^{1/2}u\|_2^2+M\bigr).
\]
Since $s>1/2$, this proves \eqref{r11tr:endpoint}. The kinetic
contribution to $J_R'$ is $2\langle u,C_Ru\rangle$.
Lemma~\ref{lem:r11tr-core} controls its error from $2mA(u)$.

For the local equation, direct integration by parts gives the nonlinear
term
\[
 -\frac{2(p-2)}p\int\operatorname{div}f_R(x-c)|u|^p.
\]
Its difference from $-4m\|u\|_p^p/p$ is bounded by
$C\int_{|x-c|>4R}|u|^p$. Critical interpolation gives
\begin{equation}\label{r11tr:local-tail}
 \int_{|x-c|>4R}|u|^p
 \le\|\eta_Ru\|_p^p\le C M^{m/N}T_R^c.
\end{equation}

For the Hartree equation, write $\rho=|u|^2$. Symmetrization gives
the nonlinear term
\[
 \iint\rho(x)\rho(y)
  [f_R(x-c)-f_R(y-c)]\cdot\nabla K(x-y)\,dx\,dy.
\]
Its difference from $-mH(u)$ vanishes if both points lie in
$B(c,4R)$. The Lipschitz bound on $f_R$ and
\eqref{r11tr:kernel} dominate the error by the absolute Riesz
interaction with at least one point outside that ball. For
$q=4N/(2N-m)$, Hardy--Littlewood--Sobolev and interpolation yield
\begin{align}
 |\text{Hartree error}|
 &\le C\|\eta_Ru\|_q^2\|u\|_q^2
 \le C_{M,B}(T_R^c)^{1/2},\label{r11tr:hartree-tail}\\
 \|w\|_q^2&\le C\|w\|_2\||D|^sw\|_2.\notag
\end{align}
Symmetrization supplies one factor $|x-y|$ at the diagonal, so HLS
also justifies passage from a regularized kernel.

Motion of the center adds
\[
 -2c'(t)\cdot\int Df_R(x-c)\operatorname{Im}(\overline u\nabla u).
\]
The difference from $-2c'\cdot P$ uses $g_R=Df_R-I$, which
vanishes on $B(c,4R)$. Its current form on $u$ is exactly its
current form on $w=\eta_Ru$: the cutoff is one on the support
of $g_R$, and its derivative contributes only a real multiple of
$|u|^2$. The rescaled first-order form estimate gives
\begin{align*}
 \left|\operatorname{Im}\int g_R\overline w\nabla w\right|
 &\le C\bigl(\||D|^{1/2}w\|_2^2+R^{-1}\|w\|_2^2\bigr)\\
 &\le C_M(T_R^c)^{1/m}+C_MR^{-1}.
\end{align*}
Here $g_R$ is a constant matrix plus a rescaled smooth compactly
supported matrix; its constant part is the ordinary momentum form.
The last inequality is interpolation between orders zero and $s$.
All current identities extend by density in $H^{1/2}$.

The principal terms are $4mE-2c'\cdot P$ in both equations.
Since $T_R^c$ is uniformly bounded, reducing the powers
$1/2,1,1/m$ to $\gamma$ proves \eqref{r11tr:virial-cut}.
Lemma~\ref{lem:r11tr-tail-equivalence} gives
\[
 (T_R^c)^\gamma\le C(D_R^c)^\gamma+C_BR^{-2\delta\gamma}.
\]
Since $2\delta\gamma=\widetilde\beta\le\beta$, this proves
\eqref{r11tr:virial-literal}. Every form and nonlinear expression
used above is continuous on bounded energy-space sets. Approximation
therefore proves the integrated statements for the solution class
specified above.
\end{proof}

\subsection{Persistent exterior kinetic density and exact averaged drift}

\begin{theorem}[Transport at a distance proportional to the time window]
\label{thm:r11tr-ballistic}
Assume $E=-e<0$ and $P=0$. For every speed bound $L$, there are
$\nu,d>0$, depending only on the fixed equation, $M,B,e,L$ and
cutoffs, such that every center $c$ with speed at most $L$ satisfies
\begin{equation}\label{r11tr:ballistic}
 \liminf_{T\to\infty}\inf_{S\ge0}\frac1T
  \int_S^{S+T}\int_{|x-c(t)|>\nu T}
                \bigl||D|^{m/2}u(t,x)\bigr|^2\,dx\,dt
 \ge d.
\end{equation}
\end{theorem}
\begin{proof}
Integrate the upper half of \eqref{r11tr:virial-literal} over
$[S,S+T]$ and use \eqref{r11tr:endpoint}:
\[
 C_1\frac1T\int_S^{S+T}(D_R^c)^\gamma\,dt
 \ge4me-C_1R^{-\widetilde\beta}-2C_0R/T.
\]
Choose $\nu>0$ with $2C_0\nu\le me$ and put $R=\nu T$,
held fixed during this time integration. For large $T$, the right
side is at least $2me$. Jensen's inequality, using $\gamma\le1$,
gives the asserted lower bound with $d=(2me/C_1)^{1/\gamma}$.
The bounds are uniform in $S$ and in all centers of the given speed.
\end{proof}

\begin{remark}
The conclusion concerns the nonlocal homogeneous kinetic density when
$m$ is not an even integer. Its exterior integral is uniformly bounded,
so the lower average forces a positive fraction of times in each long
window to have nonzero exterior kinetic density. This is not a
finite-speed theorem or a finite radiation-budget estimate, and does
not exclude indefinite replenishment of exterior activity.
\end{remark}

\begin{theorem}[Exact drift under averaged kinetic localization]
\label{thm:r11tr-drift}
For either equation and a Lipschitz center $c$, suppose
\begin{equation}\label{r11tr:average-tight}
 \lim_{R\to\infty}\limsup_{T\to\infty}\frac1T
                  \int_0^T D_R^c(t)\,dt=0.
\end{equation}
Then the limit exists and
\begin{equation}\label{r11tr:drift}
 \lim_{T\to\infty}\frac{P\cdot[c(T)-c(0)]}{T}=2mE.
\end{equation}
In particular $E<0$ and $P=0$ exclude
\eqref{r11tr:average-tight}.
\end{theorem}
\begin{proof}
Integrate the two-sided estimate \eqref{r11tr:virial-literal},
divide by $T$, and use \eqref{r11tr:endpoint} and Jensen.
For fixed $R$, the limiting discrepancy from \eqref{r11tr:drift}
is at most
\[
 C\left[R^{-\widetilde\beta}+
       \left(\limsup_{T\to\infty}\frac1T
                    \int_0^TD_R^c(t)\,dt\right)^\gamma\right].
\]
Let $R\to\infty$.
\end{proof}

The drift theorem holds at every energy sign. For a traveling wave,
it reproduces $v\cdot P=2mE$. Its content for a general localized
trajectory does not require a rigid profile or periodicity.

\subsection{Uniform localization with an arbitrary center}

\begin{theorem}[No regularity assumption on the localization center]
\label{thm:r11tr-center}
Suppose $E<0$ and an arbitrary map $b:[0,\infty)\to\mathbb R^N$
satisfies
\begin{equation}\label{r11tr:uniform-tight}
 \sup_{t\ge0}D_R^b(t)\longrightarrow0\qquad(R\to\infty).
\end{equation}
Then $b$ remains within a bounded distance of a Lipschitz center,
and
\[
 \lim_{T\to\infty}\frac{P\cdot[b(T)-b(0)]}{T}=2mE.
\]
Consequently no solution obeying \eqref{r11tr:bound},
$E<0$, $P=0$, and \eqref{r11tr:uniform-tight} exists.
\end{theorem}
\begin{proof}
Lemma~\ref{lem:r11tr-tail-equivalence} first gives uniform
localization of $T_R^b$. In the local case take $z=p$; negative
energy gives $\|u(t)\|_p^p\ge p|E|$. In the Hartree case take
$z=q=4N/(2N-m)$. Since
$H(u)=2A(u)+4|E|\ge4|E|$ and
$|H(u)|\le C\|u\|_q^4$, the $L^q$ norm also has a uniform
positive lower bound. Estimates \eqref{r11tr:local-tail} and
\eqref{r11tr:hartree-tail}, applied at smaller cutoff radii,
show that the $L^z$ norm outside $B(b(t),2R)$ tends to zero
uniformly in time.

The equation gives $\|u_t\|_{H^{-s}}\le C_B$. For the local term,
this follows by H\"older and $H^s\hookrightarrow L^p$; for Hartree,
use
\[
 \left|\int(K*|u|^2)u\overline h\right|
 \le C\|u\|_q^3\|h\|_q\le C_B\|h\|_{H^s}.
\]
Choose $\sigma$ strictly between $N(1/2-1/z)$ and $s$.
Sobolev embedding and interpolation yield
\begin{equation}\label{r11tr:time-continuity}
 \|u(t)-u(t')\|_{L^z}
 \le C_B|t-t'|^{(s-\sigma)/(2s)}.
\end{equation}
Fix $R_0$ so that the exterior $L^z$ norm is much smaller than
the positive core norm, and then $\tau>0$ so that
\eqref{r11tr:time-continuity} is smaller still for
$|t-t'|\le\tau$. If $|b(t)-b(t')|>4R_0$, the two balls of
radius $2R_0$ are disjoint. Restriction to the core ball at $t'$
then contradicts \eqref{r11tr:time-continuity}. Hence
$|b(t)-b(t')|\le4R_0$ whenever $|t-t'|\le\tau$.

Linear interpolation through $b(n\tau)$ produces a center $c$
with $\|c'\|_\infty\le4R_0/\tau$ and
$\sup_t|c(t)-b(t)|\le8R_0$. A bounded displacement preserves
literal kinetic-tail localization by inclusion of balls. Apply
Theorem~\ref{thm:r11tr-drift}; the bounded displacement also
does not affect the drift limit.
\end{proof}

The localization assumption allows a diffuse fixed-mass tail whose
homogeneous kinetic norm tends to zero. Thus it is weaker than
precompactness in $H^s$ modulo translations. No uniform control of
high-frequency tails is needed.

\subsection{A stronger exterior-gradient theorem for quartic Hartree}

For $a(\xi)=|\xi|^4$, $N>4$, and a kernel satisfying
\eqref{r11tr:kernel} with $m=4$, put
\[
 G_R^c(t)=\int_{|x-c(t)|>R}|\nabla u(t,x)|^2\,dx.
\]
Use the usual radial cutoff $f_R=\nabla\phi_R$ with
$D^2\phi_R\le I$, equal to $x$ on the core. The biharmonic
kinetic term is bounded above by
$8\|\Delta u\|_2^2+O_B(R^{-2})$. In the Hartree estimate,
$q=2N/(N-2)$ and homogeneous Sobolev give
\[
 |\text{Hartree error}|
 \le C\|\eta_Ru\|_q^2\|u\|_q^2
 \le C_B\|\nabla(\eta_Ru)\|_2^2
 \le C_B(G_R^c+R^{-2}).
\]
The center-current error is at most
$C\sqrt M L(G_R^c)^{1/2}$. Consequently
\begin{equation}\label{r11tr:quartic-gradient}
 J_R'\le16E-2c'\cdot P+C\bigl((G_R^c)^{1/2}+R^{-2}\bigr).
\end{equation}

\begin{corollary}[Quartic nonlocal exterior-gradient persistence]
For the quartic Hartree equation just specified, a forward-global
uniformly $H^2$-bounded solution with $E<0$ and $P=0$ satisfies
\[
 \liminf_{T\to\infty}\inf_{S\ge0}\frac1T
 \int_S^{S+T}G_{\nu T}^c(t)\,dt\ge d>0
\]
for every bounded-speed center, with constants depending on its speed
bound. Uniform gradient localization about an arbitrary center is
impossible under those assumptions.
\end{corollary}
\begin{proof}
The time-window integration of \eqref{r11tr:quartic-gradient}
is identical to the proof of Theorem~\ref{thm:r11tr-ballistic}.
For an arbitrary uniformly localizing center, the $L^q$ lower bound,
the Sobolev estimate on the exterior gradient, and the temporal
continuity argument in Theorem~\ref{thm:r11tr-center} produce a
bounded-speed center. Fixed-radius integration then contradicts
$E<0,P=0$.
\end{proof}

This gradient conclusion uses the isotropic fourth-order kinetic sign;
it is not claimed for arbitrary anisotropic $a$. The general two-sided
theorem instead gives exact drift at the full kinetic-localization
level.

\subsection{Relation to earlier virial arguments}
Localized virial and compact-orbit rigidity are established methods.
Boulenger--Himmelsbach--Lenzmann, \emph{Blowup for fractional NLS},
arXiv:1509.08845, develop the fractional localized virial framework.
Dinh, arXiv:1808.07368, Proposition~1.5, already proves finite-time
blowup or an unbounded kinetic subsequence for nonradial mass-critical
isotropic dispersion of degree $1<m<2$ with negative energy, without
a localization hypothesis. That specialization is therefore not a
new open-problem resolution here. Cho--Hwang--Kwon--Lee,
arXiv:1208.2302, give earlier negative-energy Hartree blowup results
under their stated assumptions.

The extension established in this section is the common anisotropic
all-degree transport argument, including nonlocal kernels, the exact
averaged drift law and its kinetic-localization form. No priority claim
is based solely on this comparison. Nor do these results exclude a
bounded global orbit that indefinitely sustains the exterior activity
required by Theorem~\ref{thm:r11tr-ballistic}.


\section{Relation to prior work and the remaining dynamical problem}
The variational stability argument follows Cazenave--Lions~\cite{CL}.
Nonhomogeneous sharp interpolation and constrained standing waves for
mixed dispersion have been studied, for example, by
Fern\'andez--Jeanjean--Mandel--Mari\c{s}~\cite{FJMM}.
The specific object computed here is the limiting concentration energy
for a bounded, possibly angular, high-frequency deficit, and the
all-data critical-mass bound below it. A proof of this statement does
not by itself establish its first priority.

Charge-response degeneracy, neutral Jordan chains, and Hamiltonian
spectral bifurcation are part of established solitary-wave theory;
see Grillakis--Shatah--Strauss~\cite{GSS},
Comech--Pelinovsky~\cite{CP}, and Lin--Zeng~\cite{LZ}.
The explicit small-velocity coefficient here includes the projection of
the scaling derivative away from translations. Its real-profile
nondegeneracy hypothesis remains a separate model-dependent question.

Bugiera--Lenzmann--Schikorra--Sok~\cite{BLSS} treat boosted ground
waves for general elliptic dispersion; their existence class includes
the critical higher-order equation considered here. Their existence
and symmetry conclusions do not supply the differentiable nondegenerate
branch assumed in our spectral calculation. In the pure quartic critical
case, Bonheure--Cast\'eras--Gou--Jeanjean~\cite{BCGJinst}, Theorem 1.2,
prove instability of radial standing ground waves by finite- or
infinite-time blowup. Our nonzero-velocity spectral asymptotic addresses
a different statement and does not replace that nonlinear result.

Sharp vector interpolation is also established; Ma--Zhao~\cite{MZ}
derive such inequalities for coupled systems.
Cheng--Guo--Hwang--Yoon~\cite{CGHY} study scattering for finite and
infinite coupled cubic systems. Our common scalar dispersion and
total-density interaction give a full internal unitary symmetry.
The polarization and compactness arguments depend on that symmetry;
they make no claim about arbitrary matrix couplings or distinct internal
frequencies.

Localized virial arguments are central in the work of
Pausader and Shao~\cite{PS}, Boulenger, Himmelsbach and Lenzmann~\cite{BHL},
and the mass-critical Hartree literature~\cite{CHKL}.
In particular, Dinh's Proposition 1.5~\cite{Dinh} already excludes a
global uniformly bounded negative-energy orbit for the isotropic
fractional mass-critical equation of order $1<m<2$, without a radial
or localization assumption. The fractional specialization of our
transport theorem is therefore not a new resolution of that question.
The precise common transport identity derived here allows arbitrary
elliptic homogeneous anisotropy, all $m>1$, and the stated Hartree class.

The principal unresolved issue along the original homogeneous quartic
line is still the near-threshold evolution of nonradial negative-energy
data. Our profile gap and transport theorem do not exclude every
persistent splitting or replenished radiation mechanism. Excluding
bounded global evolution would also leave the separate possibility of
infinite-time growth. Likewise, identifying a positive real spectral
value does not supply a nonlinear instability proof for every symbol
and dimension, and a static scale curvature does not by itself identify
a Hamiltonian oscillation frequency. These are precise missing steps
between the proved results and a complete classification of the flow.

\clearpage
\begin{thebibliography}{14}
\raggedright
\bibitem{companion} Research continuation,
\emph{Thresholds, spectral geometry, and persistent radiation in
mass-critical dispersive equations}, 8 October 2026.
Companion manuscript and complete LaTeX sources included in the source
package; in particular the threshold compactification and modified
dispersion variational sections.
\bibitem{CL} T. Cazenave and P.-L. Lions,
\emph{Orbital stability of standing waves for some nonlinear Schr\"odinger equations},
Comm. Math. Phys. \textbf{85} (1982), 549--561.
\href{https://doi.org/10.1007/BF01403504}{doi:10.1007/BF01403504}.
\bibitem{FJMM} A. Fern\'andez, L. Jeanjean, R. Mandel and M. Mari\c{s},
\emph{Non-homogeneous Gagliardo--Nirenberg inequalities in
$\mathbb R^N$ and application to a biharmonic non-linear Schr\"odinger equation},
\href{https://arxiv.org/abs/2010.01448}{arXiv:2010.01448}.
\bibitem{GSS} M. Grillakis, J. Shatah and W. Strauss,
\emph{Stability theory of solitary waves in the presence of symmetry, II},
J. Funct. Anal. \textbf{94} (1990), 308--348.
\href{https://doi.org/10.1016/0022-1236(90)90016-E}{doi:10.1016/0022-1236(90)90016-E}.
\bibitem{CP} A. Comech and D. Pelinovsky,
\emph{Purely nonlinear instability of standing waves with minimal energy},
Comm. Pure Appl. Math. \textbf{56} (2003), 1565--1607.
\href{https://doi.org/10.1002/cpa.10104}{doi:10.1002/cpa.10104}.
\bibitem{LZ} Z. Lin and C. Zeng,
\emph{Instability, index theorem, and exponential trichotomy for linear
Hamiltonian PDEs},
\href{https://arxiv.org/abs/1703.04016}{arXiv:1703.04016}.
\bibitem{BLSS} L. Bugiera, E. Lenzmann, A. Schikorra and J. Sok,
\emph{On symmetry of traveling solitary waves for dispersion generalized NLS},
\href{https://arxiv.org/abs/1908.06771}{arXiv:1908.06771}.
\bibitem{BCGJinst} D. Bonheure, J.-B. Cast\'eras, T. Gou and L. Jeanjean,
\emph{Strong instability of ground states for the fourth order
Schr\"odinger equation},
\href{https://arxiv.org/abs/1703.07977}{arXiv:1703.07977}.
\bibitem{MZ} L. Ma and L. Zhao,
\emph{Uniqueness of ground states of some coupled nonlinear Schr\"odinger
systems and their application},
\href{https://arxiv.org/abs/0708.0284}{arXiv:0708.0284}.
\bibitem{CGHY} X. Cheng, Z. Guo, G. Hwang and H. Yoon,
\emph{Global well-posedness and scattering of the two dimensional cubic
focusing nonlinear Schr\"odinger system},
\href{https://arxiv.org/abs/2202.10757}{arXiv:2202.10757}.
\bibitem{PS} B. Pausader and S. Shao,
\emph{The mass-critical fourth-order Schr\"odinger equation in high dimensions},
\href{https://arxiv.org/abs/0906.1547}{arXiv:0906.1547}.
\bibitem{BHL} T. Boulenger, D. Himmelsbach and E. Lenzmann,
\emph{Blowup for fractional NLS},
\href{https://arxiv.org/abs/1509.08845}{arXiv:1509.08845}.
\bibitem{Dinh} V. D. Dinh,
\emph{Blow-up criteria for fractional nonlinear Schr\"odinger equations},
\href{https://arxiv.org/abs/1808.07368}{arXiv:1808.07368v1}.
\bibitem{CHKL} Y. Cho, G. Hwang, S. Kwon and S. Lee,
\emph{On finite time blowup for mass-critical Hartree equations},
\href{https://arxiv.org/abs/1208.2302}{arXiv:1208.2302}.
\end{thebibliography}
\end{document}
